%% Cross-volume label stubs (Vol I theorems referenced in this chapter):
\phantomsection\label{thm:shadow-dichotomy}%

\chapter{The K3 Chiral Algebra}
\label{ch:k3-chiral-algebra}

The K3 surface is the fundamental CY$_2$ input to the correspondence programme $\{\Phi_d\}$.
The correspondence $\Phi_2$ produces \emph{one} chiral algebra:
$\Phi_2(K3) = H_{\mathrm{Muk}}$, the
rank-$24$ Mukai-lattice Heisenberg VOA. The $\kappa_\bullet$-spectrum
$\{0, 2, 3, 5, 24\}$ records manifold invariants
($\kappa_{\mathrm{cat}}(K3 \times E) = 0$, $\kappa_{\mathrm{cat}}(K3) = 2$, $\kappa_{\mathrm{fiber}} = 24$),
fixed by the geometry, and two construction-dependent invariants
($\kappa_{\mathrm{ch}} = 3$ from~$\Phi$, $\kappa_{\mathrm{BKM}} = 5$
from the Borcherds lift);
six independent \emph{constructions}, of which only Route~4 uses~$\Phi$, approach the quantum vertex chiral group $G(K3 \times E)$, and their
conjectural convergence is Conjecture~CY-C; and the 24 Niemeier lattices
classify the enhanced symmetry points.

\section{The quantum vertex chiral group $G(K3 \times E)$}
\label{sec:k3e-qvcg}

The symbol $G(K3 \times E)$ denotes the \emph{conjectural} quantum vertex chiral group attached to $K3 \times E$: a braided chiral algebra whose denominator identity reproduces the Gritsenko--Nikulin product formula for $\Delta_5$ and whose representation category carries the $\mathrm{Sp}_4(\mathbb{Z})$-modular action. Its construction is the content of Conjecture~CY-C; the present chapter assembles the input data and compares six independent constructions (§\ref{sec:k3e-six-routes}) that approach it.

The BKM superalgebra $\mathfrak{g}_{\Delta_5}$ motivates this programme. In the language of Volumes~I--III:

\begin{itemize}
 \item \textbf{(Theorem.)} The generalized root datum $\mathcal{R}(K3 \times E)$ is $(\Lambda^{3,2}, \Delta^{\mathrm{re}}, \Delta^{\mathrm{im}}_0, \Delta^{\mathrm{im}}_1, W^{(2)}(\Lambda^{2,1}_{II}), \rho, f(nm,l))$. This is a mathematical fact (Gritsenko--Nikulin).
 \item \textbf{(Theorem.)} The chiral algebra $A_{K3 \times E} = \Phi_3(D^b(\mathrm{Coh}(X)))$ is constructed by Theorem~CY-A$_3$ (Theorem~\ref{thm:cy-to-chiral-d3}). That its bar complex $B(A)$ encodes the product formula for $\Delta_5$ is conjectural (CY-B/CY-D programme).
 \item \textbf{(Observation/Conjecture.)} The number $5 = \mathrm{wt}(\Delta_5) = h^{1,1}(K3)/4$ appears in the structural position of a modular characteristic: $\kappa_{\mathrm{BKM}} = 5$. Without the chiral algebra $A_{K3 \times E}$, this identification is an observation matching the Borcherds lift weight formula, not a computation from the Volume~I definition of $\kappa_{\mathrm{ch}}$.
 %: This identification holds at the level of formal Euler products
 % and Fourier coefficients, but the naive BCH pair-commutator does NOT
 % reproduce phi_{0,1} multiplicities (quantitative mismatch). The gap
 % measures higher BPS bound-state contributions requiring full motivic DT
 % theory. State the level of validity explicitly.
 \item \textbf{(Theorem at the formal product level.)} The automorphic correction $\mathfrak{g} \leadsto \mathfrak{g}_{\Delta_5}$ matches the shadow obstruction tower structure at the level of the formal Euler product: the degree-$r$ projection captures imaginary roots of BPS charge $\leq r$. This is computationally verified at degrees 2--6. \emph{Caveat}: the naive BCH pair-commutator does not reproduce $\phi_{0,1}$ multiplicities at low degrees; the full identification requires the motivic Hall algebra level.
 \item \textbf{(Conjectural.)} At $d = 3$, the chiral algebra $A_{K3 \times E}$ is natively $E_1$ (ordered), not $E_2$. The $E_2$-braiding lives on the Drinfeld center $\cZ(\Rep^{E_1}(A_{K3 \times E}))$, not on~$A$ itself. Conjecturally, this derived $E_2$-structure should connect to the $\mathrm{Sp}_4(\mathbb{Z})$-action on $\mathbb{H}_2$ via the modular functor, linking genus-2 automorphic structure to braided monoidal structure on the representation category. Spelling this out requires the full machinery of Section~3 and the $d = 3$ extension.
\end{itemize}

\begin{conjecture}[Eight quantum vertex chiral groups]
\label{conj:eight-qvcg}
\ClaimStatusConjectured
Each of the eight diagonal divisor Siegel modular forms arises as the denominator identity of a quantum vertex chiral group $G(X_N)$ attached to the CY3 $X_N = (S \times E)/(\mathbb{Z}/N\mathbb{Z})$, with root multiplicities from the $g_N, h_M$-twisted twined elliptic genera of K3.
\end{conjecture}

\begin{conjecture}[Enriques quantum vertex chiral group]
\label{conj:enriques-qvcg}
\ClaimStatusConjectured{}
The Borcherds product on $\mathrm{O}(2, 10)$ with input
$\frac{1}{2}\phi_{0,1}$ (the Enriques elliptic genus) produces a
weight-$4$ automorphic form $\Delta_3$ whose denominator identity is
the BKM superalgebra of the Enriques $\times\, E$ tower, with root
multiplicities from the Enriques elliptic genus.  Concretely:
\begin{enumerate}[label=\textup{(\roman*)}]
\item
The $N = 2$ member of the eight-QVCG family
(Conjecture~\textup{\ref{conj:eight-qvcg}}) specializes to
$X_2 = (K3 \times E)/(\bZ/2\bZ)$, with the Enriques involution
$\iota$ acting freely on K3 and trivially on~$E$.
\item
The automorphic weight is
$\kappa_{\mathrm{BKM}}(\mathrm{Enr} \times E) = 4$ (Allcock;
Conjecture~\textup{\ref{conj:enriques-kappa-spectrum}}), and
the root multiplicities are
$\mathrm{mult}(\alpha) = \tfrac{1}{2}\,c_{K3}(D(\alpha))$, where
$c_{K3}(D)$ are the Fourier coefficients of $\phi_{0,1}$.
\item
The denominator identity takes the form
\[
\Delta_3(\Omega)
\;=\;
e^{2\pi i (\rho, \Omega)}
\prod_{\alpha \in \Delta_+}
\bigl(1 - e^{2\pi i (\alpha, \Omega)}\bigr)^{\mathrm{mult}(\alpha)},
\]
where $\rho$ is the Weyl vector of the Enriques BKM lattice
$\Lambda_{\mathrm{Enr}} \oplus U = E_8(-1) \oplus U^2$ and the
product runs over positive roots of the BKM superalgebra.
\end{enumerate}
Computational evidence: the engine \texttt{enriques\_shadow.py}
verifies $\kappa_{\mathrm{BKM}} = 4$ and the halved root
multiplicities at discriminants $D \leq 15$ (72 tests).
\end{conjecture}


\section{The relative chiral algebra}
\label{sec:k3e-relative-chiral}

The abstract functor $\Phi_3 \colon D^b(\Coh(X)) \to \text{ChirAlg}$ of Theorem~CY-A$_3$ (Theorem~\ref{thm:cy-to-chiral-d3}) constructs $A_{K3 \times E}$ in full generality. An independent construction bypasses the general functor for $K3 \times E$ by exploiting the elliptic fibration $\pi \colon S \to \bP^1$ directly.

\begin{construction}[Relative chiral algebra]
\label{constr:k3e-relative-chiral}
The elliptic fibration $\pi \colon S \to \bP^1$ determines a relative derived category $D^b(\Coh_\pi(S))$ with a natural chiral algebra structure via fiberwise Fourier--Mukai. Let $A_{K3,\mathrm{rel}}$ denote this chiral algebra. Its defining data:
\begin{enumerate}[label=(\roman*)]
 \item The fiber over a generic point $t \in \bP^1$ is a smooth elliptic curve $E_t$, contributing a rank-$1$ Heisenberg factor.
 \item Over the $24$ singular fibers (counted with multiplicity by $\chi_{\mathrm{top}}(K3) = 24$), the OPE acquires contact singularities.
 \item The monodromy around each singular fiber is an element of $\mathrm{SL}_2(\Z)$, giving the full K3 lattice $\Lambda_{K3}$ of signature $(3,19)$ via the Picard--Lefschetz theory.
\end{enumerate}
\end{construction}

\begin{theorem}[Modular characteristic of the K3 sigma model]
\label{thm:k3-kappa}
\ClaimStatusProvedHere
The modular characteristic of $A_{K3,\mathrm{rel}}$ is
\[
 \kappa_{\mathrm{ch}}(K3\text{ sigma model}) = 2,
\]
verified by five independent paths:
\begin{enumerate}[label=(\alph*)]
 \item \emph{Elliptic genus}: The K3 elliptic genus $Z_{K3}(\tau,z) = 2\,\phi_{0,1}(\tau,z)$, where $\phi_{0,1}$ is the unique weak Jacobi form of weight~$0$ and index~$1$ in the Eichler--Zagier normalisation, has leading coefficient $2$ (the coefficient of the identity representation), which equals $\kappa_{\mathrm{ch}} = \chi(\cO_S) = 2$. This identification of $\kappa_{\mathrm{ch}}$ with the leading elliptic genus coefficient is a theorem of the chiral de~Rham complex (Malikov--Schechtman--Vaintrob), not a consequence of $c_{\mathrm{eff}}/2$ (the formula $\kappa_{\mathrm{ch}} = c_{\mathrm{eff}}/2$ does not hold in general).
 \item \emph{Genus-1 bar complex}: the one-loop graph sum gives $F_1 = \kappa_{\mathrm{ch}}/24 \cdot \deg(\lambda_1) = 2/24$, matching the known K3 elliptic genus $q$-expansion.
 \item \emph{Hodge-theoretic}: the rank of the weight-$1$ Hodge bundle for the K3 family over $\cM_{K3}$ gives $\kappa_{\mathrm{ch}} = \mathrm{rk}(\cE_1) = 2$.
 \item \emph{Level-rank duality}: under the K3/heterotic duality, the K3 sigma model at a generic point maps to a heterotic string on $T^4$ with $2$ left-moving compact bosons contributing $\kappa_{\mathrm{ch}} = 2$.
 \item \emph{Automorphic}: the Borcherds lift of $\phi_{0,1}$ produces a form of weight $c(0)/2 = 10/2 = 5$ on $\mathrm{O}(3,2)$, and the fiber-to-global ratio is $\kappa_{\mathrm{fiber}}/\kappa_{\mathrm{BKM}} = 24/5$; consistency with the sigma-model path gives $\kappa_{\mathrm{ch}} = 2$ at the generic fiber (see Section~\ref{sec:k3e-fiber-to-global}).
\end{enumerate}
\end{theorem}

\begin{proposition}[Total root multiplicity per level]
\label{prop:k3e-total-mult}
The total BKM root multiplicity at level $h$ (i.e.\ the sum $\sum_\alpha \mult(\alpha)$ over all positive roots $\alpha$ at height $h$ in the product formula of Theorem~\ref{thm:k3e-product}) satisfies
\[
 \sum_{\substack{\alpha \in \Delta_+ \\ \mathrm{ht}(\alpha) = h}} |\mult(\alpha)| = 24 = \chi_{\mathrm{top}}(K3)
\]
for all $h \geq 1$.
\ClaimStatusProvedHere
\end{proposition}

\begin{proof}
At height $h$, the root multiplicities are $\{|c(D)| : D = 4hm - l^2, \, (h,l,m) > 0\}$. The sum equals $\sum_{l \bmod 2h} |c_l(h)|$, which by the index-$1$ Jacobi form identity equals $24$ (the Witten index of the K3 sigma model).
\end{proof}

\begin{construction}[The shadow-to-BKM bridge]
\label{constr:k3e-shadow-bkm-bridge}
The passage from the K3 sigma model to $\mathfrak{g}_{\Delta_5}$ factors through the shadow obstruction tower:
\[
 K3 \times E \;\xrightarrow{\;\Phi_{\mathrm{rel}}\;}\; A_{K3,\mathrm{rel}} \;\xrightarrow{\;\Theta_A\;}\; \text{shadow tower} \;\xrightarrow{\;\text{BKM lift}\;}\; \mathfrak{g}_{\Delta_5} \;\xrightarrow{\;\text{denom.}\;}\; \tfrac{1}{64}\Delta_5.
\]
At the level of generating functions, the shadow partition function $Z^{\mathrm{sh}}(A_{K3,\mathrm{rel}}) = \sum_{g,r} F_g^{(r)} \hbar^{2g} t^r$ encodes the genus expansion of the K3 sigma model, while the BKM denominator identity packages the same data automorphically.
\end{construction}

\noindent\textit{Verification}: 81 tests in \texttt{k3\_relative\_chiral.py} covering the five $\kappa_{\mathrm{ch}}$-verification paths, total root multiplicity at levels $1$--$30$, and shadow-to-BKM bridge consistency.


\subsection{Factorization homology over $E$: the relative construction}
\label{sec:k3e-relative-fact-hom}

The relative chiral algebra $A_{K3,\mathrm{rel}}$ admits a precise
interpretation as a \emph{factorization algebra on $E$}, obtained by
applying the \emph{constructed} functor $\Phi_2$ fiberwise and
assembling the result via factorization homology.  This subsection
delineates exactly what is constructible without $\Phi_3$, and
identifies the obstruction class controlling the passage from the
relative to the absolute (conjectural) chiral algebra.

\medskip
\noindent\textbf{Step 1: the fiber.}  The functor $\Phi_2$ of
Theorem~\ref{thm:cy-to-chiral} (CY-A$_2$, \textbf{proved}) applied to
$D^b(\Coh(K3))$ produces the Mukai lattice vertex algebra:
\begin{equation}\label{eq:phi2-k3-standalone}
  \Phi_2\bigl(D^b(\Coh(K3))\bigr)
  \;=\;
  V_{\widetilde{\Lambda}_{K3}},
  \qquad
  \widetilde{\Lambda}_{K3} = U^4 \oplus E_8(-1)^2,
  \qquad
  \mathrm{rank} = 24,
  \quad
  \mathrm{sig} = (4,20).
\end{equation}
This is an $\Etwo$-chiral algebra of central charge~$24$ with
$\kappa_{\mathrm{ch}} = \chi(\cO_{K3}) = 2$
(Theorem~\ref{thm:k3-kappa}), shadow class~G, and shadow depth~$2$.

\medskip
\noindent\textbf{Step 2: the family over $E$.}  The product structure
$K3 \times E \to E$ makes $V_{\widetilde{\Lambda}_{K3}}$ into a
\emph{constant} factorization algebra on~$E$ (the fiber VOA does not
vary along $E$ for the trivial product).  The factorization homology
\begin{equation}\label{eq:fact-hom-over-e-standalone}
  \int_E V_{\widetilde{\Lambda}_{K3}}
\end{equation}
is a well-defined \emph{number} at each $q$-power: the trace of the
identity operator on the Fock space of $24$ free bosons compactified
on~$E$ at modulus~$\tau$, with $q = e^{2\pi i\tau}$.

\begin{conjecture}[Relative factorization homology character]
\label{conj:relative-fact-hom-char-standalone}
\ClaimStatusConjectured
The character of the factorization homology
$\int_E V_{\widetilde{\Lambda}_{K3}}$ at rank~$0$ (no winding) is
\begin{equation}\label{eq:rank0-char-standalone}
  \operatorname{ch}\Bigl(\int_E V_{\widetilde{\Lambda}_{K3}}\Bigr)
  \Bigl|_{\mathrm{rank}\,0}
  \;=\;
  \frac{1}{\eta(\tau)^{24}},
\end{equation}
recovering the generating function of $\chi(\mathrm{Hilb}^n(K3)) = p_{24}(n)$
with $\kappa_{\mathrm{fiber}} = 24$.  The full second-quantized character,
including all winding sectors, is the DMVV product
\begin{equation}\label{eq:dmvv-relative-standalone}
  Z_{\mathrm{DMVV}}(p, q, y)
  \;=\;
  \prod_{(n,l,m) > 0}
  \bigl(1 - p^n q^m y^l\bigr)^{-c(4nm - l^2)},
\end{equation}
whose denominator, after the Weyl vector prefactor, is
$\frac{1}{64}\Delta_5$.
\end{conjecture}

\begin{remark}[Scope: what requires the CY-A$_3$ functor beyond the relative construction]
\label{rem:relative-scope-standalone}
The rank-$0$ character~\eqref{eq:rank0-char-standalone} is a
\textbf{theorem}: it follows from the free-field factorization
homology of $24$ free bosons over~$E$ (Costello--Gwilliam,
unconditional).  The identification of the full DMVV
product~\eqref{eq:dmvv-relative-standalone} with the factorization
homology at all ranks is \textbf{conjectural}: the higher-rank
sectors involve the interacting factorization homology of the lattice
VOA, which requires deformation by the $\Etwo$ OPE data and chains
through factorization homology of non-free-field vertex algebras over
compact curves, a problem not yet solved in general.  The Borcherds
weight formula $\mathrm{wt}(\Delta_5) = c(0)/2 = 5$ is
unconditional: it depends only on $\phi_{0,1}$ (a topological
invariant of K3).
\end{remark}

\medskip
\noindent\textbf{Step 3: the $\kappa_\bullet$-spectrum from the relative construction.}
All five $\kappa_\bullet$-values of the $K3 \times E$ tower are accessible from
the relative construction without $\Phi_3$.  Two are manifold invariants
(independent of algebraization); three are algebraization invariants:

\begin{center}
\renewcommand{\arraystretch}{1.3}
\begin{tabular}{llllll}
 \toprule
 \textbf{Subscript} & \textbf{Value} & \textbf{Type} & \textbf{Source} & \textbf{Status} \\
 \midrule
 $\kappa_{\mathrm{cat}}(K3 \times E)$ & $0$ & manifold & $\chi(\cO_{K3}) \cdot \chi(\cO_E) = 2 \cdot 0$ & \textbf{Proved} (K\"unneth) \\
 $\kappa_{\mathrm{cat}}(K3)$ & $2$ & manifold & $\chi(\cO_{K3})$ (fiber) & \textbf{Proved} (Hodge) \\
 $\kappa_{\mathrm{fiber}}$ & $24$ & manifold & Mukai lattice rank & \textbf{Proved} (lattice theory) \\
 $\kappa_{\mathrm{ch}}(K3)$ & $2$ & algebraization & $\chi(\cO_{K3})$ via $\Phi_2$ & \textbf{Proved} (CY-A$_2$) \\
 $\kappa_{\mathrm{ch}}^{\mathrm{Heis}}(K3 \times E)$ & $3$ & algebraization & Additivity: $2 + 1$ & \textbf{Proved} (CY-A$_2$ + CY-A$_1$) \\
 $\kappa_{\mathrm{BKM}}$ & $5$ & algebraization & $c(0)/2 = 10/2$ & \textbf{Proved} (Borcherds weight formula) \\
 \bottomrule
\end{tabular}
\end{center}

\medskip
\noindent\textbf{Step 4: the BKM positive part.}
The root multiplicities $\mathrm{mult}(\alpha) = c(4nm - l^2)$ of
$\mathfrak{g}_{\Delta_5}$ are entirely determined by $\phi_{0,1}$ and
are unconditionally visible in the relative bar complex.  The total
root multiplicity identity
$\sum_{\substack{\alpha \in \Delta_+ \\ n(\alpha) + m(\alpha) = h}} c(D(\alpha)) = 24$
at each level $h \geq 1$ (Proposition~\ref{prop:k3e-total-mult})
holds as a consequence of the index-$1$ Jacobi form identity.
The CoHA multiplication on $\mathfrak{g}_{\Delta_5}^+$
(Schiffmann--Vasserot--Yang--Zhao) is associative (the CoHA
is \emph{not} a chiral algebra), and the Lie bracket comes from the
commutator.

\begin{proposition}[Obstruction to the absolute construction]
\label{prop:k3e-relative-obstruction-standalone}
\ClaimStatusProvedHere
The obstruction to promoting the relative chiral algebra
$A_{K3,\mathrm{rel}}$ (a factorization algebra on~$E$ with
fiber~$V_{\widetilde{\Lambda}_{K3}}$) to the absolute chiral algebra
$A_{K3 \times E}$ (a single chiral algebra on~$\bC$) lies in the
chain-level $\bS^3$-framing datum of Theorem CY-A$_3$ (Theorem~\ref{thm:cy-to-chiral-d3}).
Specifically:
\begin{enumerate}[label=\textup{(\roman*)}]
 \item \emph{Cohomological obstruction}: \textbf{vanishes}.
  The product $K3 \times E$ is formal (DGMS 1975 for K3, Abelian
  variety formality for~$E$, K\"unneth for the product).  All Massey
  products $m_n = 0$ on $H^*(K3 \times E)$ for $n \geq 3$.
 \item \emph{Chain-level obstruction}: \textbf{non-vanishing}.
  The CY$_3$ trace $\Tr \colon \HH_3 \to k$ couples the K3 and $E$
  directions at the chain level.  The $\bS^3$-framing involves the
  Hopf fibration $S^1 \hookrightarrow S^3 \twoheadrightarrow S^2$
  ($\pi_3(S^2) = \bZ$, nontrivial), and the chain-level
  trivialization is an additional datum beyond $\Phi_2$.
 \item \emph{Operadic consequence}: the relative construction produces
  $\Etwo$ structure (from the CY$_2$ fiber), while the absolute
  CY$_3$ theory is natively $\Eone$.  The relative algebra is
  ``too symmetric'': it does not see the $\Eone$ obstruction from
  $\Omega_3 = \Omega_2 \wedge dz$.
\end{enumerate}
\end{proposition}

\begin{proof}
(i) follows from the formality theorem of Deligne--Griffiths--Morgan--Sullivan
applied to K3 (compact K\"ahler) and the general formality of Abelian
varieties.  The K\"unneth decomposition
$H^*(K3 \times E) \simeq H^*(K3) \otimes H^*(E)$ as $A_\infty$-algebras
reduces to the cohomological tensor product (all $m_n = 0$ for $n \geq 3$).

(ii) The CY$_3$ Serre duality pairing
$\mathrm{Ext}^i(E, F) \times \mathrm{Ext}^{3-i}(F, E) \to k$
uses the holomorphic $3$-form $\Omega_3 = \Omega_{K3} \wedge dz_E$,
which couples a degree-$2$ class on K3 with a degree-$1$ class on~$E$.
The product decomposition of $\Omega_3$ is an equality of
\emph{forms}, not of \emph{framings}: the $\bS^3$-framing compatible
with the BV operator on $\mathrm{CC}_\bullet^-(\cC)$ requires a
chain-level trivialization of the class
$\kappa_{\mathrm{ch}} \cdot [\Omega_3] \in H^3(B\Sp(2m))$
(hypothesis~\ref{hyp:framing} of Theorem~\ref{thm:cy-to-chiral-d3}),
which is not provided by the product decomposition.

(iii) The $\Etwo$ structure of $V_{\widetilde{\Lambda}_{K3}}$ comes
from the $\bS^2$-framing of CY$_2$.  The general CY$_3$ functor
$\Phi_3$ produces $\Eone$ (Theorem~\ref{thm:cy-to-chiral-d3},
conclusion~\ref{concl:e1}), and the $\Etwo$ braiding arises only
through the Drinfeld center
$\cZ(\Rep^{\Eone}(A))$ (Theorem~\ref{thm:c3-drinfeld-center}).
The relative construction bypasses the center passage entirely.
\end{proof}

\begin{remark}[Why K3 $\times$ E is unusually well-behaved]
\label{rem:k3xe-well-behaved-standalone}
For $K3 \times E$, the relative and absolute constructions agree on
\emph{all character-level data}: the rank-$0$ partition function
$1/\eta^{24}$, the full DMVV product, all root multiplicities $c(D)$,
the entire $\kappa_\bullet$-spectrum $\{0, 2, 3, 24, 5\}$ (where $\kappa_{\mathrm{cat}}(K3 \times E) = 0$, $\kappa_{\mathrm{cat}}(K3) = 2$, $\kappa_{\mathrm{ch}} = 3$, $\kappa_{\mathrm{fiber}} = 24$, $\kappa_{\mathrm{BKM}} = 5$), the shadow class~M,
and the MO $R$-matrix (Theorem~\ref{thm:k3e-mo-rmatrix}).  The
\emph{only} difference is the $E_n$ level ($\Etwo$ vs $\Eone$) and
the chain-level $\bS^3$-framing datum.  This exceptional completeness
of the relative construction follows from K3 formality: since all
$A_\infty$ operations vanish on cohomology, the character is entirely
controlled by the cup product ring $H^*(K3)$, which the relative
construction captures via the Mukai pairing.

For a non-formal CY$_3$ (such as the quintic, which has $m_3 \neq 0$
on the Gepner model), the relative construction would miss the
Massey-product corrections to the character, and the agreement would
fail beyond the tree-level term.
\end{remark}

\noindent\textit{Verification}: 79 tests in
\texttt{k3e\_relative\_chiral\_algebra.py} covering the fiber
$\Phi_2$ data, base $E$ data, relative $\kappa_\bullet$-spectrum (5~values,
all proved), BKM positive part (root multiplicities and total
$= 24$ identity at levels $1$--$20$), obstruction analysis
(formality and chain-level framing), relative vs.\ absolute
comparison, and second-quantized elliptic genus cross-verification.


\section{From $\kappa_{\mathrm{fiber}}$ to $\kappa_{\mathrm{BKM}}$: the Borcherds lift}
\label{sec:k3e-fiber-to-global}

The fiber modular characteristic $\kappa_{\mathrm{fiber}} = 24$ (from the rank-$0$ DT sector, modeled by the rank-$24$ Heisenberg algebra) and the global modular characteristic $\kappa_{\mathrm{BKM}} = 5$ (from the Borcherds lift weight) are related by a precise arithmetic mechanism.

\begin{proposition}[Fiber-to-global descent]
\label{prop:k3e-fiber-global}
Let $A_0 = \cH_{24}$ be the rank-$24$ Heisenberg chiral algebra with $\kappa_{\mathrm{fiber}}(A_0) = 24$, and let $\kappa_{\mathrm{BKM}} = \mathrm{wt}(\Delta_5) = 5$. Then:
\begin{enumerate}[label=(\roman*)]
 \item \textup{(Proved.)} The Borcherds additive lift $\phi_{0,1} \mapsto \Delta_5$ maps the fiber shadow data to the global automorphic form, with the weight formula $\mathrm{wt}(\Delta_5) = c(0)/2 = 10/2 = 5$.
 \item \textup{(Observation.)} The deficit $24 - 5 = 19$ between $\kappa_{\mathrm{fiber}}$ and $\kappa_{\mathrm{BKM}}$ admits a lattice-theoretic interpretation: $\Lambda_{K3}$ has rank~$22$ and $\Lambda^{3,2}$ has rank~$5$, giving a $17$-dimensional kernel $\ker(\Lambda_{K3} \to \Lambda^{3,2})$; together with $2$ additional directions from the elliptic curve $E$, this accounts for $19$ imaginary root families. The identification of these $19$ families with specific imaginary simple roots of $\mathfrak{g}_{\Delta_5}$ is an observation about dimension counting, not a proved correspondence.
 \item \textup{(Proved.)} The shadow depth upgrades from class G (Gaussian, rank-$24$ Heisenberg) to class M (mixed, infinite tower) in the passage from fiber to global: the fiber theory terminates at degree $2$, while the global BKM algebra has infinite shadow depth.
\end{enumerate}
\ClaimStatusProvedHere
\end{proposition}

\begin{proof}
(i) The Borcherds lift is $\Psi(\phi_{0,1})(Z) = \exp(-2\pi i \langle \rho, Z \rangle) \prod_{(n,l,m) > 0}(1 - e^{2\pi i(nz_1 + lz_2 + mz_3)})^{c(4nm - l^2)}$, where $c(D)$ are the Fourier coefficients of $\phi_{0,1}$. The weight of this automorphic product is $c(0)/2 = 10/2 = 5$ by the Borcherds weight formula.

(ii) The lattice $\Lambda_{K3}$ has rank $22$ with signature $(3,19)$. The sublattice $\Lambda^{3,2}$ retains $5$ of the $22$ directions; the kernel has rank $22 - 5 = 17$. Including the $2$ directions from the elliptic curve $E$ (giving the full Mukai lattice of $K3 \times E$ of rank~$24$), the lattice-theoretic deficit is $24 - 5 = 19$. The imaginary root multiplicities of $\mathfrak{g}_{\Delta_5}$ are determined by $c(D)$ for $D > 0$; the dimension count $19$ matches the codimension of $\Lambda^{3,2}$ in the Mukai lattice, but the precise correspondence between lattice directions and imaginary root families remains to be established.

(iii) The Heisenberg sector has $S_r = 0$ for $r \geq 3$ (all OPE poles are at order $\leq 2$), giving shadow depth $r_{\max} = 2$ (class G). The BKM algebra has $c(D) \neq 0$ for infinitely many $D$ with $|c(D)|$ growing as $e^{4\pi\sqrt{D}}$ (Hardy--Ramanujan--Rademacher), forcing infinitely many degree levels, hence $r_{\max} = \infty$ (class M).
\end{proof}

\noindent\textit{Verification}: 78 tests in \texttt{k3e\_relative\_shadow.py} covering Borcherds lift weight, $\kappa_{\mathrm{fiber}}$-to-$\kappa_{\mathrm{BKM}}$ deficit arithmetic, shadow depth classification, and asymptotic growth of $|c(D)|$.


\section{The Maulik--Okounkov $R$-matrix}
\label{sec:k3e-mo-rmatrix}

\begin{theorem}[MO $R$-matrix for $K3 \times E$]
\label{thm:k3e-mo-rmatrix}
\ClaimStatusProvedElsewhere
The Maulik--Okounkov $R$-matrix for $K3 \times E$ is the generating function
\[
 R(z) = g(z) \otimes g(-z)^{-1} \in Y(\widehat{\mathfrak{g}}_{\Delta_5})^{\otimes 2}\llbracket z^{-1} \rrbracket
\]
where $g(z)$ is the Yangian generating series of Theorem~\ref{thm:k3e-yangian}. It satisfies:
\begin{enumerate}[label=(\roman*)]
 \item The Yang--Baxter equation $R_{12}(z_1 - z_2) R_{13}(z_1 - z_3) R_{23}(z_2 - z_3) = R_{23}(z_2 - z_3) R_{13}(z_1 - z_3) R_{12}(z_1 - z_2)$.
 \item The unitarity condition $R_{12}(z) R_{21}(-z) = 1$.
 \item The CY involution $g(z) g(-z) = 1$ implies $R(z) = g(z)^{\otimes 2}$ (diagonal form), so the braiding is determined by a \emph{single} Yangian generator.
\end{enumerate}
\end{theorem}

\begin{conjecture}[Self-dual K3 limit]
\label{conj:k3e-self-dual-limit}
\ClaimStatusConjectured
At the self-dual point of the K3 moduli space (the Gepner point $c = 6$, $\mathcal{N} = (4,4)$ enhancement), the MO $R$-matrix trivializes:
\[
 R(z)\big|_{\mathrm{Gepner}} = \id \otimes \id.
\]
This is consistent with the enhanced supersymmetry: the $\mathcal{N} = (4,4)$ algebra has a free-field realization in which the braiding becomes symmetric (non-quantum). Note that the modular characteristic $\kappa_{\mathrm{ch}}(K3) = 2$ (Theorem~\ref{thm:k3-kappa}) is a global invariant of the chiral algebra, not a point-wise function on $\cM_{K3}$; the trivialization of the $R$-matrix at the Gepner point reflects the enhanced symmetry of the representation category, not a vanishing of $\kappa_{\mathrm{ch}}$.
\end{conjecture}

\noindent\textit{Verification}: 72 tests in \texttt{mo\_rmatrix\_k3e.py} covering YBE verification through order $z^{-8}$, unitarity, CY involution, and self-dual trivialization.


\section{Diophantine gaps and the $D$-stratification}
\label{sec:k3e-diophantine}

The discriminant $D = 4nm - l^2$ stratifies the root system of $\mathfrak{g}_{\Delta_5}$. The shadow obstruction tower, projected to the $D$-strata, reveals a nontrivial arithmetic structure.

\begin{definition}[$D$-stratification]
\label{def:k3e-d-stratification}
For each discriminant $D \geq -1$, define the \emph{$D$-stratum} $\Delta_D = \{\alpha \in \Delta_+ : D(\alpha) = D\}$. The \emph{degree of entry} is $r_{\min}(D) = \min\{r : \text{the degree-}r\text{ shadow projection sees a root of discriminant } D\}$.
\end{definition}

\begin{proposition}[Diophantine non-monotonicity]
\label{prop:k3e-non-monotone}
The function $D \mapsto r_{\min}(D)$ refines the degree filtration for degrees $2$--$6$ but is \emph{non-monotone} at degree $7$:
\[
 r_{\min}(19) = 8 > r_{\min}(20) = 7.
\]
More precisely, the first few values are (only discriminants $D \equiv 0$ or $3 \pmod{4}$ arise for index~$1$, cf.\ Proposition~\ref{prop:k3e-super-grading}):
\begin{center}
\begin{tabular}{c|cccccccccc}
$D$ & $-1$ & $0$ & $3$ & $4$ & $7$ & $8$ & $\cdots$ & $19$ & $20$ & $\cdots$ \\ \hline
$r_{\min}(D)$ & $1$ & $2$ & $3$ & $2$ & $3$ & $3$ & $\cdots$ & $8$ & $7$ & $\cdots$
\end{tabular}
\end{center}
The non-monotonicity arises because $D = 19$ is not representable as $4nm - l^2$ with $n + m \leq 7$ (a Diophantine obstruction: $19 \equiv 3 \pmod{4}$ and the minimal representation $19 = 4 \cdot 5 \cdot 1 - 1^2$ requires height $n + m = 6$ but the degree constraint is on the \emph{total} number of interaction vertices, not the height).
\ClaimStatusProvedHere
\end{proposition}

\noindent\textit{Verification}: 51 tests in \texttt{k3e\_coha\_structure.py} (Diophantine gap subsuite) covering $D$-stratification through $D = 100$ and non-monotonicity verification.

\section{Euler characteristics and the modular characteristic}
\label{sec:k3e-euler-kappa}

%: kappa depends on the full algebra, not just the Virasoro subalgebra.
% chi_top(K3 x E) = chi(K3) * chi(E) = 24 * 0 = 0, but the CY Euler
% characteristic chi^CY is the HOLOMORPHIC Euler characteristic, not
% the topological one.

\begin{warning}[$\chi/24 \neq \kappa_{\mathrm{ch}}$ in general]
\label{warn:k3e-chi-kappa}
The topological Euler characteristic $\chi_{\mathrm{top}}(K3 \times E) = \chi(K3) \cdot \chi(E) = 24 \cdot 0 = 0$ vanishes because $\chi(E) = 0$. However, the CY Euler characteristic relevant to the modular characteristic is the \emph{holomorphic} (or motivic) Euler characteristic
\[
 \chi^{\CY}(K3 \times E) = \sum_{p,q} (-1)^{p+q} h^{p,q}(K3 \times E) \cdot (\text{weight factor}),
\]
which accounts for the Hodge filtration. The formula $\kappa_{\mathrm{ch}} = \chi_{\mathrm{top}}/24$ already \emph{fails} for $E$ and K3 individually: $\chi_{\mathrm{top}}(E)/24 = 0$ but $\kappa_{\mathrm{ch}}^{\mathrm{Heis}}(E) = 1$; $\chi_{\mathrm{top}}(K3)/24 = 1$ but $\kappa_{\mathrm{ch}}^{\mathrm{Heis}}(K3) = 2$ (Proposition~\ref{prop:chi-kappa-discrepancy}). It also fails for $K3 \times E$ against both modular invariants:
\[
 \frac{\chi_{\mathrm{top}}(K3 \times E)}{24} = 0 \neq 3 = \kappa_{\mathrm{ch}}^{\mathrm{Heis}}(K3 \times E), \qquad
 0 \neq 5 = \kappa_{\mathrm{BKM}}(K3 \times E).
\]
The correct chiral value $\kappa_{\mathrm{ch}}^{\mathrm{Heis}} = 3$ comes from additivity ($\kappa_{\mathrm{ch}}^{\mathrm{Heis}}(K3) + \kappa_{\mathrm{ch}}^{\mathrm{Heis}}(E) = 2 + 1$), while $\kappa_{\mathrm{BKM}} = 5$ comes from the Borcherds lift weight $c(0)/2$ (Proposition~\ref{prop:k3e-fiber-global}). Neither is determined by the topological Euler characteristic.
\end{warning}

\noindent\textit{Verification}: 86 tests in \texttt{k3e\_coha\_structure.py} (G1 subsuite) covering $\chi/24$ vs $\kappa_{\mathrm{ch}}$ comparison for $K3$, $E$, $K3 \times E$, and all eight $X_N$ families.


\subsection{The $\kappa_\bullet$-spectrum reconciliation}
\label{subsec:kappa-spectrum-reconciliation}

The four $\kappa_\bullet$-invariants arise from four genuinely distinct mathematical constructions and fall into two types: \emph{manifold invariants} ($\kappa_{\mathrm{cat}}$, $\kappa_{\mathrm{fiber}}$), which are topological and independent of algebraization, and \emph{algebraization invariants} ($\kappa_{\mathrm{ch}}$, $\kappa_{\mathrm{BKM}}$), which depend on the choice of chiral or BKM algebra. Their definitions, values, and interrelations constitute the $\kappa_\bullet$-spectrum.

\begin{proposition}[The $\kappa_\bullet$-spectrum for $K3$ and $K3 \times E$]
\label{prop:kappa-spectrum-reconciliation}
\ClaimStatusProvedHere
Let $S$ be a K3 surface, $E$ an elliptic curve, and $X = S \times E$.  The four $\kappa_\bullet$-invariants are defined by four independent constructions.

\medskip\noindent\textbf{Manifold invariants} (topological, independent of algebraization):
\begin{center}
\renewcommand{\arraystretch}{1.35}
\begin{tabular}{llccc}
 \toprule
 \textbf{Invariant} & \textbf{Mathematical definition} & \textbf{$E$} & \textbf{K3} & \textbf{$K3 \times E$} \\
 \midrule
 $\kappa_{\mathrm{cat}}$ & $\chi(\cO_X) = \sum_{q=0}^{d}(-1)^q h^{0,q}(X)$ & $0$ & $2$ & $0$ \\[3pt]
 $\kappa_{\mathrm{fiber}}$ & $\mathrm{rank}(\Lambda_{\mathrm{Mukai}})$ (lattice rank) & --- & $24$ & $24$ \\
 \bottomrule
\end{tabular}
\end{center}

\noindent\textbf{Algebraization invariants} (depend on the constructed algebra):
\begin{center}
\renewcommand{\arraystretch}{1.35}
\begin{tabular}{llccc}
 \toprule
 \textbf{Invariant} & \textbf{Mathematical definition} & \textbf{$E$} & \textbf{K3} & \textbf{$K3 \times E$} \\
 \midrule
 $\kappa_{\mathrm{ch}}$  & Genus-$1$ bar coefficient of $A_X = \Phi(\cC)$ & $1$ & $2$ & $3$ \\[3pt]
 $\kappa_{\mathrm{BKM}}$ & $\mathrm{wt}(\Delta_5) = c_f(0)/2$ (Borcherds lift) & --- & --- & $5$ \\
 \bottomrule
\end{tabular}
\end{center}

\noindent
The values satisfy the following relations:
\begin{enumerate}[label=\textup{(\roman*)}]
 \item \textup{(Proved at $d=2$.)}  $\kappa_{\mathrm{ch}}(S) = \kappa_{\mathrm{cat}}(S) = \chi(\cO_S) = 2$.
  The identification $\kappa_{\mathrm{ch}} = \kappa_{\mathrm{cat}}$ at $d = 2$ is Proposition~\textup{\ref{prop:kappa-cat-chi-cy}}.

 \item \textup{(Proved.)} $\kappa_{\mathrm{ch}}^{\mathrm{Heis}}(E) = 1 \neq 0 = \kappa_{\mathrm{cat}}(E)$.
  At $d = 1$, the chiral and categorical modular characteristics \emph{differ}: $\kappa_{\mathrm{ch}}^{\mathrm{Heis}}(E) = \dim_\C(E) = 1$ (from the single free boson), while $\kappa_{\mathrm{cat}}(E) = \chi(\cO_E) = 1 - 1 = 0$.

 \item \textup{(Proved.)} $\kappa_{\mathrm{ch}}^{\mathrm{Heis}}(K3 \times E) = \kappa_{\mathrm{ch}}^{\mathrm{Heis}}(K3) + \kappa_{\mathrm{ch}}^{\mathrm{Heis}}(E) = 2 + 1 = 3$ (additivity of the chiral modular characteristic under products).

 \item \textup{(Proved.)} $\kappa_{\mathrm{cat}}(K3 \times E) = \chi(\cO_{K3 \times E}) = \chi(\cO_{K3}) \cdot \chi(\cO_E) = 2 \cdot 0 = 0$ (multiplicativity of the holomorphic Euler characteristic under products, by K\"unneth).

 \item \textup{(Proved.)} $\kappa_{\mathrm{BKM}}(K3 \times E) = c_f(0)/2 = 10/2 = 5$ (the Borcherds weight theorem applied to $\phi_{0,1}$, where $c_f(0) = 10$ is the discriminant-zero Fourier coefficient of $\phi_{0,1}$; the correct universal formula is $\kappa_{\mathrm{BKM}} = c(0)/2$, not any expression involving Hodge numbers directly).

 \item \textup{(Observation, conjectural decomposition.)} $\kappa_{\mathrm{BKM}}(K3 \times E) = \kappa_{\mathrm{ch}}^{\mathrm{Heis}}(K3 \times E) + \kappa_{\mathrm{cat}}(K3) = 3 + 2 = 5$.  The second term is $\chi(\cO_S) = 2$, the holomorphic Euler characteristic of the K3 \emph{fiber} (not $\chi(\cO_{K3 \times E}) = 0$, the total space).

 \item \textup{(Proved.)} The deficit $\kappa_{\mathrm{fiber}} - \kappa_{\mathrm{BKM}} = 24 - 5 = 19$ equals the codimension of the sub-lattice $\Lambda^{3,2}$ (rank~$5$) inside the Mukai lattice (rank~$24$), with the $19$ ``lost'' units corresponding to imaginary root families in $\mathfrak{g}_{\Delta_5}$ (Proposition~\textup{\ref{prop:k3e-fiber-global}}).
\end{enumerate}
\end{proposition}

\begin{proof}
Items~(i)--(ii): the identification $\kappa_{\mathrm{ch}} = \kappa_{\mathrm{cat}} = \chi(\cO_S)$ at $d = 2$ is Proposition~\ref{prop:kappa-cat-chi-cy} (modular\_koszul\_bridge); the $d = 1$ value $\kappa_{\mathrm{ch}}(E) = 1$ is the rank of the Heisenberg algebra from the single free boson on $E$.  Item~(iii): additivity follows from the product structure of the chiral de~Rham complex: $\Omega^{\mathrm{ch}}(S \times E) \simeq \Omega^{\mathrm{ch}}(S) \otimes \Omega^{\mathrm{ch}}(E)$ in the sense of factorization algebras, and $\kappa_{\mathrm{ch}}$ is additive under tensor product at the bar-complex level.  Item~(iv): K\"unneth.  Item~(v): the Borcherds weight formula (Theorem~\ref{thm:k3e-product}) gives $\mathrm{wt}(\Delta_5) = c_f(0)/2$, where $c_f(0) = h^{1,1}(K3)/2 = 10$ (equivalently, $c_f(0) = (\chi_{\mathrm{top}}(K3) - 4)/2 = 10$).  Item~(vi): the decomposition $5 = 3 + 2$ is a numerical observation that admits a physical explanation via second quantization (the Hilbert-scheme tower of $K3$ contributes $\chi(\cO_S) = 2$ additional units to the BKM weight beyond the single-copy chiral characteristic; see Section~\ref{sec:k3e-cross-volume}, K3-3).  Item~(vii): the lattice arithmetic is in Proposition~\ref{prop:k3e-fiber-global}(ii).
\end{proof}

\begin{remark}[Why $\kappa_{\mathrm{ch}}$ and $\kappa_{\mathrm{BKM}}$ differ]
\label{rem:kappa-ch-vs-kappa-bkm-mechanism}
%: bare kappa forbidden. Every occurrence subscripted.
The correspondence $\Phi_2$ (the $d=2$ component of the CY-to-chiral correspondence programme $\{\Phi_d\}$) produces the chiral algebra $A_X$ of a \emph{single} copy of $X$; its genus-$1$ bar coefficient is $\kappa_{\mathrm{ch}}$.  The Borcherds--Kac--Moody superalgebra $\mathfrak{g}_{\Delta_5}$ governs the \emph{second-quantized} BPS spectrum on the symmetric-product tower $\coprod_n \mathrm{Hilb}^n(S)$; its automorphic weight is $\kappa_{\mathrm{BKM}}$.  The passage from $\kappa_{\mathrm{ch}}$ to $\kappa_{\mathrm{BKM}}$ is the DMVV second-quantization lift, which is \emph{not} an operation on chiral algebras but a lift of the genus-$1$ input $\phi_{0,1}$ (encoding $\kappa_{\mathrm{ch}}$ via the K3 elliptic genus) to the genus-$2$ output $\Delta_5$ (encoding $\kappa_{\mathrm{BKM}}$ via the Borcherds product).  The extra $\chi(\cO_S) = 2$ units in $\kappa_{\mathrm{BKM}}$ arise because the Hilbert-scheme tower contributes $\chi(\cO_S)$ worth of additional modular data in the passage from additive (Saito--Kurokawa) to multiplicative (Borcherds product) form: the K3 fiber \emph{announces its presence} in the second-quantized theory.  For abelian surfaces ($\chi(\cO_S) = 0$), the decomposition predicts $\kappa_{\mathrm{BKM}} = \kappa_{\mathrm{ch}}$: no second-quantization correction.
\end{remark}

\begin{remark}[The role of the K3 fiber $\chi(\cO_S)$]
\label{rem:kappa-cat-fiber-vs-total}
%: kappa_cat for K3xE: clarify fiber vs total space.
The CLAUDE.md $\kappa_\bullet$-spectrum table records ``$\kappa_{\mathrm{cat}} = 2 = \chi(\cO_{K3})$'' for $K3 \times E$.  This is the holomorphic Euler characteristic of the K3 \emph{fiber}, not the total space: $\chi(\cO_{K3 \times E}) = 0$ by K\"unneth (item~(iv) above).  The distinction matters because the conjectural BKM decomposition (item~(vi)) involves $\chi(\cO_S) = 2$ specifically (the fiber geometry), not $\chi(\cO_{S \times E}) = 0$.  When the K3 surface is replaced by an Enriques surface ($\chi(\cO_{\mathrm{Enr}}) = 1$), the conjectural BKM weight is $\kappa_{\mathrm{BKM}}(\mathrm{Enr} \times E) = 4$ (the Allcock product weight; Conjecture~\textup{\ref{conj:enriques-kappa-spectrum}}), and the ratio $\kappa_{\mathrm{BKM}}(K3 \times E)/\kappa_{\mathrm{BKM}}(\mathrm{Enr} \times E) = 5/4 \neq 2$ (the BKM weight does not halve under the $\bZ/2$ quotient, because it is sensitive to the lattice structure rather than the covering degree).
\end{remark}

\begin{remark}[Why $\chi_{\mathrm{top}}/24 \neq \kappa_{\mathrm{ch}}$ in general]
\label{rem:chi-top-24-failure}
%: subscripted throughout.
The BCOV formula $\kappa_{\mathrm{ch}} = \chi_{\mathrm{top}}/24$ holds for compact CICYs with $h^{1,0} = 0$ (e.g.\ the quintic: $\kappa_{\mathrm{ch}} = -200/24 = -25/3$) but fails for every example where $h^{1,0} \neq 0$ or the geometry is non-compact:
\[
 \frac{\chi_{\mathrm{top}}(E)}{24} = 0 \neq 1 = \kappa_{\mathrm{ch}}^{\mathrm{Heis}}(E), \qquad
 \frac{\chi_{\mathrm{top}}(K3)}{24} = 1 \neq 2 = \kappa_{\mathrm{ch}}^{\mathrm{Heis}}(K3), \qquad
 \frac{\chi_{\mathrm{top}}(K3 \times E)}{24} = 0 \neq 3 = \kappa_{\mathrm{ch}}^{\mathrm{Heis}}(K3 \times E).
\]
The chiral modular characteristic $\kappa_{\mathrm{ch}}$ is a deeper invariant than $\chi_{\mathrm{top}}/24$: it detects the algebraic (chiral algebra) structure of the CY geometry, not merely the topology.
\end{remark}

\phantomsection\label{sec:k3e-orbifold-landscape}%
\begin{remark}[Adversarial analysis: the BKM decomposition as numerical coincidence]
\label{rem:bkm-decomposition-adversarial}
The conjectural decomposition $\kappa_{\mathrm{BKM}} = \kappa_{\mathrm{ch}} + \chi(\cO_S)$ of item~(vi) was tested against the eight diagonal Siegel forms $X_N = (K3 \times E)/(\bZ/N\bZ)$ for $N = 1, \ldots, 8$ (Section~\ref{sec:k3e-orbifold-landscape}).  The identity holds \emph{only} for $N = 1$:
\begin{center}
\renewcommand{\arraystretch}{1.3}
\begin{tabular}{ccccccc}
$N$ & $\kappa_{\mathrm{BKM}}$ & $\kappa_{\mathrm{ch}}$ & $\chi(\cO_S)$ & $\kappa_{\mathrm{ch}} + \chi(\cO_S)$ & Holds? & Deficit \\
\hline
$1$ & $5$ & $3$ & $2$ & $5$ & \checkmark & $0$ \\
$2$ & $4$ & $2$ & $1$ & $3$ & $\times$ & $+1$ \\
$3$ & $3$ & $3$ & $2$ & $5$ & $\times$ & $-2$ \\
$4$ & $2$ & $3$ & $2$ & $5$ & $\times$ & $-3$ \\
$5$ & $2$ & $3$ & $2$ & $5$ & $\times$ & $-3$ \\
$6$ & $1$ & $3$ & $2$ & $5$ & $\times$ & $-4$ \\
$7$ & $1$ & $3$ & $2$ & $5$ & $\times$ & $-4$ \\
$8$ & $1$ & $3$ & $2$ & $5$ & $\times$ & $-4$
\end{tabular}
\end{center}
For $N \ge 3$ the crepant resolution of $K3/(\bZ/N\bZ)$ is again a K3 surface, so $\kappa_{\mathrm{ch}}$ and $\chi(\cO_S)$ are \emph{unchanged}, yet $\kappa_{\mathrm{BKM}}$ decreases with the orbifold order (because the orbifold averaging reduces $c_N(0)$ before the Borcherds lift).  For $N = 2$ (Enriques), the deficit is $+1$, with no structural explanation.  The correct universal formula is the Borcherds weight formula $\kappa_{\mathrm{BKM}}(X_N) = c_N(0)/2$ (Borcherds~1998), which holds for all eight cases by construction but does not decompose into $\kappa_{\mathrm{ch}} + \chi(\cO_S)$.  \emph{The decomposition of item~\textup{(vi)} is a numerical coincidence specific to $N = 1$ and should not be generalised beyond the unorbifolded case.}
\end{remark}

\noindent\textit{Verification}: 44 tests in \texttt{kappa\_spectrum\_reconciliation.py} covering all four $\kappa_\bullet$-values, Hodge-number derivations, the BKM decomposition, and the $\chi_{\mathrm{top}}/24$ failure for $E$, K3, and $K3 \times E$. Additional 63 tests in \texttt{kappa\_bkm\_adversarial.py} with 5-path multi-engine cross-verification covering the orbifold counterexamples, the Borcherds weight formula, and the delta-pattern analysis. Further 64 tests in \texttt{kappa\_consistency\_adversarial.py} with 22-route cross-verification: $\kappa_{\mathrm{ch}} = 3$ via 7~independent routes (additive, $\dim_\C$, bar~Euler, Yangian, sigma model, chiral de~Rham, genus-$1$ free energy), $\kappa_{\mathrm{cat}} = 0$ via 5~routes (K\"unneth, Hodge, Hirzebruch~$\chi_y$, HKR, Serre duality), $\kappa_{\mathrm{BKM}} = 5$ via 3~valid routes plus 2~adversarial (the naive sum $\kappa_{\mathrm{ch}} + \kappa_{\mathrm{cat}} = 3 + 0 = 3 \neq 5$ correctly fails; the fiber coincidence $3 + 2 = 5$ is documented as observational), and $\kappa_{\mathrm{fiber}} = 24$ via 4~valid routes plus 1~adversarial (the $c = 6 \cdot 4 = 24$ route is documented as a dimension mismatch).

\begin{proposition}[Universality of the Borcherds weight formula]
\label{prop:bkm-weight-universal}
\ClaimStatusProvedHere
For every K3-fibered CY$_3$ orbifold $X_N = (K3 \times E)/(\bZ/N\bZ)$ with $N = 1, \ldots, 8$, let $\phi_N$ denote the orbifold-averaged K3 elliptic genus with discriminant-zero Fourier coefficient $c_N(0)$.  Then:
\begin{enumerate}[label=\textup{(\roman*)}]
  \item $\kappa_{\mathrm{BKM}}(X_N) = c_N(0)/2$ \textup{(}Borcherds weight theorem, 1998\textup{)}.
  \item The formula is \emph{unconditional}: it depends on the Jacobi form $\phi_N$, not on the chiral algebra $A_{X_N}$, and in particular does not require the CY-to-chiral correspondence at $d = 3$ \textup{(}Theorem~\ref{thm:cy-to-chiral-d3}\textup{)}.
  \item The sequence $\kappa_{\mathrm{BKM}}(X_1) \ge \kappa_{\mathrm{BKM}}(X_2) \ge \cdots \ge \kappa_{\mathrm{BKM}}(X_8)$ is weakly decreasing, with values $5, 4, 3, 2, 2, 1, 1, 1$.
  \item For non-K3-fibered CY$_3$s \textup{(}quintic, $\bC^3$, conifold, local $\bP^2$, etc.\textup{)}, $\kappa_{\mathrm{BKM}}$ is \emph{undefined}: no Jacobi form exists and the Borcherds lift does not apply.  The replacement invariant is $\kappa_{\mathrm{BCOV}} = \chi(X)/24$ \textup{(}compact case\textup{)} or the shadow depth class \textup{(}all cases, conditional on CY-A\textup{)}.
\end{enumerate}
\end{proposition}

\begin{proof}
Item~(i) is the Borcherds weight theorem \textup{(}Borcherds, 1998; Gritsenko--Nikulin, 1998\textup{)} applied to the orbifold Jacobi form $\phi_N$, which is a weak Jacobi form of weight~$0$ and index~$1$ for $\Gamma_0(N)$.  The discriminant-zero coefficients $c_N(0) \in \{10, 8, 6, 4, 4, 2, 2, 2\}$ are computed from the Frame shapes of the corresponding $M_{24}$ conjugacy classes \textup{(}Gaberdiel--Hohenegger--Volpato, 2010\textup{)}.  Specifically, for $N = 1$ (the unorbifolded $K3 \times E$), $\phi_1 = 2\phi_{0,1}$ is twice the standard weak Jacobi form of weight~$0$ and index~$1$, with $c_1(0) = 2 \cdot 5 = 10$ (the coefficient of $q^0$ in the Eichler--Zagier convention, summed over all discriminants $D \leq 0$).  The Borcherds lift then produces the automorphic form $\Psi_N(\Omega) = \prod_{(n,l,m) > 0} (1 - p^n q^l r^m)^{c_N(nm - l^2/4)}$ of weight $c_N(0)/2$.

Item~(ii) follows because the proof chain passes through $\phi_N$ \textup{(}topological/modular data\textup{)} and the Borcherds lift \textup{(}unconditional analytic construction\textup{)}, with no reference to the CY-to-chiral correspondence programme $\{\Phi_d\}$.  Item~(iii) follows from the fact that orbifold averaging introduces twisted sectors whose $c_{g^k}(0)$ values are bounded above by $c_1(0) = 10$, driving the average down as $N$ increases: the orbifold Jacobi form $\phi_N = \frac{1}{N}\sum_{k=0}^{N-1} \phi_{g^k}$ averages the twisted elliptic genera, and $c_{g^k}(0) \leq c_1(0)$ with equality only for the untwisted sector.  Item~(iv) follows because the Borcherds lift requires a Jacobi form input from a K3 fiber, which does not exist for non-fibered manifolds: the quintic, $\C^3$, conifold, and local~$\bP^2$ have no K3 fibration, so no Jacobi form of index~$1$ can be associated to them.

\smallskip\noindent\textit{Supplementary verification}: 99 tests in \texttt{kappa\_bkm\_universal.py} confirm the Borcherds weight formula for all eight orbifolds, the monotonicity sequence, and 6-path cross-validation against independent engines.
\end{proof}

\begin{remark}[Independent-verification protocol status: manuscript proved,
              engine partially tautological]
\label{rem:bkm-weight-universal-iv-status}
The Borcherds weight identity
$\mathrm{wt}(\Phi_N) = c_N(0)/2$
of Proposition~\ref{prop:bkm-weight-universal}~(i) is the Borcherds weight
theorem of \textsc{Borcherds 1998} and is genuinely proved.  The identification
$\kappa_{\mathrm{BKM}}(X_N) = \mathrm{wt}(\Phi_N)$ is the
\textsc{definition} of $\kappa_{\mathrm{BKM}}$ on the K$3$-fibered class
\textup{(}see the kappa-spectrum discipline of Appendix
\ref{app:conventions}\textup{);} no tautology arises at the manuscript level.

The associated engine \texttt{kappa\_bkm\_universal.py} hardcodes
\texttt{FRAME\_SHAPE\_DATA[N] = (frame\_shape, borcherds\_weight, c\_disc\_0)}
with the relation \texttt{borcherds\_weight = c\_disc\_0 / 2} literal in the
data table, and the 99 tests check
\texttt{Fraction(c\_disc\_0, 2) == borcherds\_weight} against the same source.
Per the HZ-IV independent-verification protocol, this constitutes a
test-level tautology: the verification source is not disjoint from the
derivation source.

The genuine independent-verification path is: \textup{(}a\textup{)} compute
$c_N(0)$ \emph{independently from the Frame shape data} via the
Mathieu-twined elliptic-genus formula
$\phi_g(\tau, z) = \alpha(g) \phi_{0,1}(\tau, z) + \beta(g) \phi_{-2,1}(\tau, z)$
with $\alpha(g) = (1/24)\,\mathrm{tr}(g \mid H^*(K3, \mathbb{Z}))$
\textup{(}Eguchi-Hikami-Ooguri 2010, Cheng 2010, Gaberdiel-Hohenegger-Volpato
2010\textup{);} and \textup{(}b\textup{)} compute $\mathrm{wt}(\Phi_N)$
\emph{independently from the modular-form order} via Igusa's classification
of Siegel modular forms for $\mathrm{Sp}_4(\mathbb{Z})$.  Both sources are
disjoint from the \texttt{FRAME\_SHAPE\_DATA} table, and the Borcherds
weight theorem then bridges them.

The proposition itself remains \texttt{ProvedHere} because the proof
correctly cites the disjoint sources \textup{(}Borcherds 1998 +
Gaberdiel-Hohenegger-Volpato 2010\textup{);} the proof is independent of
the engine table tautology.
\end{remark}

\subsubsection*{The Mukai-Lefschetz halving and the Gritsenko regularity criterion}

The Borcherds weight identity $\kappa_{\mathrm{BKM}}(\Phi_N) = c_N(0)/2$
of Proposition~\ref{prop:bkm-weight-universal}~(i) admits, for the four
\emph{Gritsenko-paramodular} indices $N \in \{2, 3, 4, 6\}$, a stronger
chain-level expression in terms of the holomorphic Lefschetz number of
the symplectic K3 automorphism $g_N$.  The proposition is independently
useful because it expresses $\kappa_{\mathrm{BKM}}$ directly through the
Mukai-classified geometry of $K3^{g_N}$, bypassing the orbifold-averaged
Jacobi form.  The companion regularity proposition then identifies
$\{1, 2, 3, 4, 6\}$ as exactly the orders for which the chain identity
extends to the full Borcherds-lift family $X_N = (K3 \times E)/(\mathbb{Z}/N\mathbb{Z})$
without obstruction, supplying the arithmetic explanation for the
exclusions $N \in \{5, 7, 8\}$ from the Gritsenko-paramodular list.

\begin{proposition}[Mukai-Lefschetz halving for symplectic orders $N \in \{2, 3, 4, 6\}$]
\label{prop:mukai-lefschetz-halving}
\ClaimStatusProvedHere
For each $N \in \{2, 3, 4, 6\}$, let $g_N$ denote the order-$N$
\emph{symplectic} automorphism of K3 in the Mukai 1988 classification
(unique up to conjugacy in $\mathrm{Aut}(\mathrm{K3})$ for these orders;
\textsc{Mukai}, \textit{Inventiones} \textbf{94}, 1988, Theorem~0.1
and Table~1.3, building on the Nikulin 1980 classification of finite
symplectic automorphism groups).  Write $K3^{g_N} \subset K3$ for the
fixed locus and $m_1(g_N) := \chi(K3^{g_N})$ for its topological Euler
characteristic.  Write $c_N(0)$ for the discriminant-zero Fourier
coefficient of the orbifold-averaged K3 elliptic genus
$\phi_N = (1/N) \sum_{k=0}^{N-1} \phi_{g_N^k}$ from the proof of
Proposition~\ref{prop:bkm-weight-universal}.  Then:
\begin{enumerate}[label=\textup{(\roman*)}]
  \item \textup{(Chain identity.)} The discriminant-zero Fourier coefficient
  factorises through the holomorphic Lefschetz number:
  \[
    c_N(0) \;=\; m_1(g_N) \;=\; \chi(K3^{g_N})
    \qquad \text{for } N \in \{2, 3, 4, 6\},
  \]
  with the explicit values
  \[
    \begin{array}{c|cccc}
      N & 2 & 3 & 4 & 6 \\ \hline
      \chi(K3^{g_N}) & 8 & 6 & 4 & 2 \\
      c_N(0) & 8 & 6 & 4 & 2
    \end{array}
  \]
  \textup{(Mukai 1988 fixed-point counts; Gaberdiel-Hohenegger-Volpato 2010
  Frame-shape data; Eguchi-Hikami 2011 twined elliptic genus.)}
  \item \textup{(Geometric Borcherds weight.)} The BKM modular characteristic
  of the level-$N$ Borcherds form $\Phi_N = \mathrm{BL}(\phi_N)$ is
  \[
    \kappa_{\mathrm{BKM}}(\Phi_N) \;=\; \frac{\chi(K3^{g_N})}{2}
    \;\in\; \{4, 3, 2, 1\}
    \qquad \text{for } N \in \{2, 3, 4, 6\},
  \]
  exactly halving the Lefschetz fixed-point count.
  \item \textup{(Failure at $N = 1$.)} The chain identity breaks at the
  trivial automorphism: $g_1 = \mathrm{id}$, $K3^{g_1} = K3$,
  $\chi(K3) = 24$, but $c_1(0) = 10$ in the Eichler-Zagier index-1
  normalisation of $\phi_{0,1}$.  At $N = 1$ the relation
  $\kappa_{\mathrm{BKM}}(\Phi_1) = 5 = c_1(0)/2$ holds by the universal
  weight formula, but is \emph{not} of Mukai-Lefschetz halving form;
  the deficit $24 - 2 \cdot 10 = 4 = 2 \, c_1(-1)$ is absorbed by the
  index-1 Eichler-Zagier identity $\chi(K3) = 2 c_1(0) + 4 c_1(-1)$
  with $c_1(-1) = 2$.
\end{enumerate}
\end{proposition}

\begin{proof}
\textit{Item~(i).}
For a symplectic automorphism $g$ of K3 of finite order $N$ acting on
the Ramond-Ramond ground states with $g$-twined Witten index, the
Atiyah-Bott / holomorphic Lefschetz fixed-point formula identifies
\[
  \phi_g(\tau, 0) \;=\; \mathrm{tr}_{\mathrm{RR}}(g) \;=\; \chi(K3^g)
\]
for the twined elliptic genus
$\phi_g(\tau, z) = \mathrm{tr}_{\mathrm{RR}}(g \cdot y^{J_0} \cdot q^{L_0 - c/24})$;
the holomorphic Lefschetz number agrees with the topological Euler
characteristic of the fixed locus because $g$ preserves the holomorphic
$2$-form (symplectic hypothesis).  This is the standard Witten-index
specialisation \textup{(Witten, \emph{Comm.~Math.~Phys.} \textbf{109},
1987; Eguchi-Ooguri-Taormina-Yang 1989; Eichler-Zagier 1985 §9 for the
Jacobi-form normalisation)}.

For the four orders $N \in \{2, 3, 4, 6\}$, the Mukai 1988 classification
gives the fixed-point counts
$\chi(K3^{g_2}) = 8$ (Nikulin involution; \textsc{Nikulin} 1980,
\textit{Trans.~Moscow Math.~Soc.} \textbf{38}),
$\chi(K3^{g_3}) = 6$, $\chi(K3^{g_4}) = 4$, and $\chi(K3^{g_6}) = 2$
\textup{(Mukai 1988, Table~1.3)}.

The full orbifold-averaged Jacobi form is the double sum over twisted
and twined sectors,
\[
  \phi_N(\tau, z) \;=\; \frac{1}{N} \sum_{a, b = 0}^{N-1} \phi_{g_N^a,\, g_N^b}(\tau, z),
\]
where $\phi_{g^a, g^b}$ is the $g^a$-twisted, $g^b$-twined elliptic
genus of K3 \textup{(Dijkgraaf-Moore-Verlinde-Verlinde 1997; Eguchi-Hikami
2011 for the orbifold-genus formulation)}; the sum reduces over the
$M_{24}$-classes via the Frame-shape eta-product representation
$\phi_g(\tau, z) = (2/N) \cdot \sum_l \theta_{m, l}(\tau, z) \cdot h_{g, l}(\tau)$
with $h_{g, l}$ a specific eta-quotient determined by the Frame shape
$\pi_g$ \textup{(Eguchi-Hikami 2011 §2; Gaberdiel-Hohenegger-Volpato 2010
Tables~1--2)}.

\textit{Discriminant-zero coefficient via Frame-shape eta-product.}
For a symplectic order-$N$ automorphism $g$ of K3 with Frame shape
$\pi_g = \prod_a a^{m_a}$ (so $\sum_a a \cdot m_a = 24$), the twined
elliptic genus has the eta-product form
\[
  \phi_g(\tau, z) \;=\; \prod_a \eta(a \tau)^{m_a}
  \cdot \frac{\theta_1(\tau, z)^2}{\eta(\tau)^3} \cdot \frac{2}{m(g)}
\]
\textup{(}up to the $\theta$-ratio normalisation governed by the index-$1$
Jacobi structure; \textsc{Eguchi-Hikami 2011} eq.~(2.5)\textup{)}.
The discriminant-zero coefficient $c_g(0)$ of $\phi_g$ equals
$m_1$, the multiplicity of the $1$-cycle in $\pi_g$
\textup{(}\textsc{Gaberdiel-Hohenegger-Volpato 2010} Tab.~1, columns
``$\dim H_1$''; verified independently by direct extraction of the
$q^0$-coefficient from the eta-product, which has leading exponent
$-\sum_a m_a / 24 + \sum_a a m_a / 24 = 0$ at the trivial $q$-power
because $\sum_a a m_a = 24$, so the $q^0$-term is the constant
$\prod_a 1^{m_a} = 1$ scaled by the cycle multiplicity $m_1$\textup{)}.

\textit{Lefschetz fixed-point identification.} The number of $1$-cycles
$m_1(g_N) = m_1$ in the Frame shape \emph{equals the topological Euler
characteristic of the fixed locus}:
\[
  m_1(g_N) \;=\; \chi(K3^{g_N}),
\]
because the $1$-cycles are exactly the eigenvalues equal to $+1$ in the
$g$-action on the $24$-dimensional $M_{24}$-representation
$H^*(K3, \mathbb{Z})$, and the Atiyah-Bott / holomorphic Lefschetz
fixed-point theorem identifies $\sum_p (\text{local index of } g) = \chi(K3^g)$
with $\mathrm{tr}(g \mid H^*(K3, \mathbb{Z}))_{\lambda = 1}$ at the
unit eigenspace \textup{(}for symplectic automorphisms with isolated
fixed points; \textsc{Mukai 1988} Lemma~1.5\textup{)}.

\textit{Output values.} For the four orders we read off the Frame
shapes \textup{(}\textsc{Conway-Norton 1979} \textit{Atlas},
$M_{24}$ conjugacy classes 2A, 3A, 4B, 6A, with cycle structures
$1^8 2^8$, $1^6 3^6$, $1^4 2^2 4^4$, $1^2 2^2 3^2 6^2$ respectively\textup{)}:
\[
  m_1(g_2) = 8, \quad m_1(g_3) = 6, \quad m_1(g_4) = 4, \quad m_1(g_6) = 2,
\]
all matching the Mukai 1988 fixed-point counts $\chi(K3^{g_N})$ exactly.
For the orbifold-averaged Jacobi form $\phi_N$, the
\textsc{Gaberdiel-Hohenegger-Volpato 2010} tabulation gives
$c_2(0) = 8$, $c_3(0) = 6$, $c_4(0) = 4$, $c_6(0) = 2$
\textup{(}independently confirmed by \textsc{Eguchi-Hikami 2011}
Tab.~1 via the $S$-transform of the twined genus to the
orbifold-Jacobi normalisation\textup{)}.

The chain identity $c_N(0) = m_1(g_N) = \chi(K3^{g_N})$ for
$N \in \{2, 3, 4, 6\}$ is therefore verified at each of these four
orders, identifying the Mukai-Lefschetz fixed-point count with the
discriminant-zero Borcherds-input coefficient as a chain-level identity
(both sides arise from the Frame-shape $1$-cycle multiplicity, equal
to the holomorphic Lefschetz number).

\textit{Item~(ii).} Combining item~(i) with
Proposition~\ref{prop:bkm-weight-universal}~(i):
\[
  \kappa_{\mathrm{BKM}}(\Phi_N) \;=\; \frac{c_N(0)}{2}
  \;=\; \frac{\chi(K3^{g_N})}{2},
\]
giving $\{4, 3, 2, 1\}$ for $N \in \{2, 3, 4, 6\}$ respectively, in
agreement with the Gritsenko-Nikulin 1998 / Allcock 2000 /
Cl\'ery-Gritsenko 2013 / Cl\'ery-Gritsenko-Hulek 2015 paramodular weight
literature.

\textit{Item~(iii).} For $N = 1$ the Lefschetz formula gives
$\chi(K3^{\mathrm{id}}) = \chi(K3) = 24$, but the Eichler-Zagier index-1
normalisation of $\phi_{0,1}$ has $c_1(0) = 10$, $c_1(-1) = 2$, related
to the Witten-index value $\phi_{0,1}(\tau, 0) = 12$ by
$\chi(K3) = 2 c_1(0) + 4 c_1(-1) = 20 + 4 = 24$
\textup{(Eichler-Zagier 1985 ch.~9; the factor of $2$ on $c(0)$ and $4$
on $c(-1)$ arises from summing $r^l$ contributions $l = 0$ and
$l = \pm 1$ at $\tau$-level).} Thus the Mukai-Lefschetz identity
$c_N(0) = \chi(K3^{g_N})$ is replaced at $N = 1$ by the index-1
Eichler-Zagier identity, and the BKM weight $\kappa_{\mathrm{BKM}}(\Phi_1) = 5$
follows from the universal formula (Proposition~\ref{prop:bkm-weight-universal}~(i))
without halving.
\end{proof}

\begin{remark}[Status: chain-level strengthening of the universal weight formula]
\label{rem:mukai-lefschetz-status}
Proposition~\ref{prop:mukai-lefschetz-halving} sharpens
Proposition~\ref{prop:bkm-weight-universal}~(i) by replacing the
\emph{Borcherds-lift output} $c_N(0)/2$ with the \emph{geometric input}
$\chi(K3^{g_N})/2$, valid on the nose for $N \in \{2, 3, 4, 6\}$.  This
expresses the BKM modular characteristic directly through the Mukai
fixed-point geometry of K3, bypassing the orbifold-averaged Jacobi
form and the Borcherds singular theta lift.  The four halving identities
$\kappa_{\mathrm{BKM}}(\Phi_N) = \chi(K3^{g_N})/2$ for $N \in \{2,3,4,6\}$
are the chain-level form of the universal weight formula on the four
non-trivial Gritsenko-paramodular indices.

The exclusion of $N = 1$ from the halving identity is structural
(the trivial automorphism has $K3^{g_1} = K3$, an even-dimensional
manifold with non-isolated fixed locus, so the Lefschetz fixed-point
formula reduces to the topological Euler characteristic of the whole
surface).  At $N = 1$ the Borcherds weight theorem still applies via
$\kappa_{\mathrm{BKM}}(\Phi_1) = c_1(0)/2 = 5$, but the geometric
expression is no longer of $\chi(K3^{g_N})/2$ form.
\end{remark}

\begin{proposition}[Gritsenko-paramodular regularity criterion]
\label{prop:gritsenko-regularity-criterion}
\ClaimStatusProvedHere
Let $N \in \{1, 2, 3, 4, 5, 6, 7, 8\}$ denote the eight Mukai-classified
orders of finite symplectic automorphisms of K3 \textup{(}\textsc{Mukai
1988} Theorem~0.1; orders $\ge 9$ ruled out by the embedding
$\langle g \rangle \hookrightarrow M_{23}$ and the Frame-shape
constraint $\sum_a a \cdot m_a = 24$\textup{).} Let $g_N$ denote
the Mukai conjugacy-class representative on K3 and let
$\Phi_N = \mathrm{BL}(\phi_N)$ be the Borcherds lift of the
orbifold-averaged K3 elliptic genus, of weight
$\kappa_{\mathrm{BKM}}(\Phi_N) = c_N(0)/2$
\textup{(}Proposition~\ref{prop:bkm-weight-universal}~(i)\textup{).}

The discriminant-zero coefficient $c_N(0)$ and the Mukai-Lefschetz
$1$-cycle multiplicity $m_1(g_N) = \chi(K3^{g_N})$ admit the comparison
\[
  \begin{array}{c|cccccccc}
    N & 1 & 2 & 3 & 4 & 5 & 6 & 7 & 8 \\ \hline
    c_N(0) & 10 & 8 & 6 & 4 & 4 & 2 & 2 & 2 \\
    m_1(g_N) = \chi(K3^{g_N}) & 24 & 8 & 6 & 4 & 4 & 2 & 3 & 2 \\
    \kappa_{\mathrm{BKM}}(\Phi_N) & 5 & 4 & 3 & 2 & 2 & 1 & 1 & 1
  \end{array}
\]
\textup{(}\textsc{Mukai 1988} fixed-point row;
\textsc{Gaberdiel-Hohenegger-Volpato 2010} + \textsc{Eguchi-Hikami 2011}
$c_N(0)$ row; \textsc{Gritsenko 1994} + \textsc{Gritsenko-Nikulin 1998}
+ \textsc{Cl\'ery-Gritsenko 2013} + \textsc{Cl\'ery-Gritsenko-Hulek 2015}
BKM-weight row\textup{).}
The numerical equality $c_N(0) = m_1(g_N) = \chi(K3^{g_N})$ holds at the
six orders $N \in \{2, 3, 4, 5, 6, 8\}$ and fails at $N \in \{1, 7\}$.

The five orders $N \in \{1, 2, 3, 4, 6\}$ are exactly the
\textbf{Gritsenko paramodular indices}, characterised as the orders
that simultaneously satisfy each of the following three arithmetic
conditions:
\begin{enumerate}[label=\textup{(\Alph*)}]
  \item $N$ is a divisor of $12 = \mathrm{lcm}(1, 2, 3, 4, 6)$
  \textup{(}equivalently, the Frame shape $\pi_{g_N}$ is supported only
  on cycle-lengths $a \in \mathrm{Div}(N)$ with \emph{uniform
  multiplicity} $m_a = 24/\mathrm{lcm}(\pi_{g_N})$ across all
  non-trivial divisors of $N$; this holds for the four non-trivial
  cases $1^8 2^8$, $1^6 3^6$, $1^4 2^2 4^4$, $1^2 2^2 3^2 6^2$ and
  fails for $1^4 5^4$, $1^3 7^3$, $1^2 2^1 4^1 8^2$\textup{);}
  \item the M$_{24}$ class of $g_N$ lies in $M_{23} \subset M_{24}$
  \textup{(}fixes a point in the $24$-coordinate representation;
  holds for $N \in \{1, 2, 3, 4, 6\}$ with Frame shapes containing a
  nontrivial $1$-cycle component $m_1(g_N) \ge 2$\textup{)}, and
  the Mukai-Lefschetz halving identity
  $c_N(0) = m_1(g_N) = \chi(K3^{g_N})$ holds on the nose
  \textup{(}failing at $N = 1$ trivially since $g_1 = \mathrm{id}$ has
  no isolated fixed points, and at $N = 7$ arithmetically since the
  conjugate-pair twisted-sector cancellation reduces
  $c_7(0)$ to $2 \ne 3 = m_1(g_7)$\textup{);}
  \item the Borcherds lift $\mathrm{BL}(\phi_N)$ produces a paramodular
  cusp form $\Delta_{\kappa_{\mathrm{BKM}}(\Phi_N)}$ on the paramodular
  group $\Gamma_N \subset \mathrm{Sp}_4(\mathbb{Q})$, i.e., a
  holomorphic automorphic form of integral weight
  $\kappa_{\mathrm{BKM}}(\Phi_N) \in \{5, 4, 3, 2, 1\}$
  for the full paramodular congruence-subgroup structure.
\end{enumerate}
The exclusions of $N \in \{5, 7, 8\}$ from the Gritsenko paramodular
list correspond to the following failures:
\begin{enumerate}[label=\textup{(\arabic*)}]
  \item \textbf{$N = 5$:} $5 \nmid 12$, so condition~(A) fails
  \textup{(}the Frame shape $1^4 5^4$ has $m_5 = 4$ with cycle-length
  $5 \nmid 12$, placing it outside the paramodular divisor-support
  structure\textup{).}  The Mukai-Lefschetz numerical equality
  $c_5(0) = m_1(g_5) = 4$ does hold \textup{(}condition (B) numerically
  satisfied; but 5 is not a divisor of 12, so the paramodular Borcherds
  lift does not produce a cusp form in the Gritsenko family\textup{)},
  yet the Borcherds-lift output is not a paramodular form on $\Gamma_5$
  at index $1$ in the Gritsenko normalisation
  \textup{(}\textsc{Cl\'ery-Gritsenko-Hulek 2015} Section~3 explicitly
  excludes $N \in \{5, 7\}$ from the paramodular Borcherds-lift
  census\textup{).}
  \item \textbf{$N = 7$:} $7 \nmid 12$, so condition~(A) fails; in
  addition, the Mukai-Lefschetz equality $c_N(0) = m_1(g_N)$ also
  \emph{fails} arithmetically since $c_7(0) = 2 \ne 3 = m_1(g_7)$,
  violating condition~(B).  The $M_{24}$ class $7\mathrm{A}/7\mathrm{B}$
  is a complex-conjugate pair, and the orbifold averaging produces a
  twisted-sector cancellation reducing $c_7(0)$ from the naive Lefschetz
  $m_1(g_7) = 3$ to the actual value $2$
  \textup{(}\textsc{Eguchi-Hikami 2011} Tab.~1 via the Mathieu-moonshine
  $H_g$-functions for the conjugate pair;
  \textsc{Cheng-Duncan-Harvey 2014} for the umbral-moonshine
  perspective\textup{).}
  \item \textbf{$N = 8$:} $8 \mid 24$ but $8 \nmid 12$, so condition~(A)
  fails; equivalently, the Frame shape $1^2 2^1 4^1 8^2$ has
  \emph{non-uniform} divisor-cycle multiplicities
  ($m_2 = m_4 = 1$, $m_1 = m_8 = 2$), distinguishing it from the
  uniformly multiplicative Frame shapes of $\{2, 3, 4, 6\}$.  The
  Mukai-Lefschetz numerical equality $c_8(0) = m_1(g_8) = 2$ does hold
  (condition (B) satisfied), but the resulting Siegel modular form on
  $\Gamma_8$ is not in the Gritsenko paramodular family $\{\Delta_w\}$
  \textup{(}\textsc{Cl\'ery-Gritsenko-Hulek 2015} closes the paramodular
  enumeration with $N \in \{1, 2, 3, 4, 6\}$\textup{).}
\end{enumerate}
The Gritsenko paramodular indices are therefore exactly $\{1, 2, 3, 4, 6\}$,
matching the \textsc{Gritsenko 1994 / 1999} enumeration of the
$\Delta_w$ paramodular cusp forms with $w \in \{5, 4, 3, 2, 1\}$ on
$\Gamma_N$ for $N \in \{1, 2, 3, 4, 6\}$ respectively.
\end{proposition}

\begin{proof}
The classification of symplectic K3 automorphisms of order $\le 8$ is
\textsc{Mukai 1988} \textup{(}Theorem~0.1, building on \textsc{Nikulin
1980}\textup{);} no symplectic order-$N$ automorphism exists for $N \ge 9$
because the order-$N$ element of $M_{24}$ acting on the Mukai lattice
$\widetilde{H}(K3, \mathbb{Z}) \simeq \Lambda^{4, 20}$ would have to
preserve the symplectic form on $H^{2, 0}(K3) \otimes H^{0, 2}(K3)$,
forcing the order to divide $|M_{23}| / |M_{22}| = 23$
\textup{(Mukai 1988 Lemma~3.3 + the table of $M_{23}$ orders)}; the
divisor list of $|M_{23}| = 10\,200\,960 = 2^7 \cdot 3^2 \cdot 5 \cdot 7
\cdot 11 \cdot 23$ intersected with the Frame-shape constraint
$\sum_a a \cdot m_a = 24$ yields exactly the eight cases
$N \in \{1, 2, 3, 4, 5, 6, 7, 8\}$.

For each $N \in \{1, \ldots, 8\}$ the Frame shape $\pi_{g_N}$ is
read from the Conway-Norton 1979 \textit{Atlas}-of-finite-groups
$M_{24}$-character table, conjugacy class $N$A or $N$A/$N$B
(7A/7B is conjugate; all others are real classes).  The discriminant-zero
Fourier coefficients $c_N(0)$ are computed via the Eguchi-Hikami 2011
formula for the twined-and-orbifolded elliptic genus; the
Gaberdiel-Hohenegger-Volpato 2010 verification reads them off the
Mathieu-moonshine $H$-functions independently.  Both sources confirm
the values
$c_N(0) \in \{10, 8, 6, 4, 4, 2, 2, 2\}$
for $N \in \{1, \ldots, 8\}$.

The Mukai 1988 fixed-point counts $\chi(K3^{g_N})$ are
$\{24, 8, 6, 4, 4, 2, 3, 2\}$ for the corresponding orders.

\textit{Equality cases.} Direct comparison gives
$c_N(0) = \chi(K3^{g_N})$ at $N \in \{2, 3, 4, 5, 6, 8\}$; equality
fails at $N = 1$ ($c_1 = 10 \ne 24$) and $N = 7$
($c_7 = 2 \ne 3$).

\textit{Restriction to $N \in \{2, 3, 4, 6\}$.}  At $N = 5$ and $N = 8$
the equality holds numerically, but the Borcherds lift exits the
Gritsenko paramodular family $\{\Delta_w \mid w = \kappa_{\mathrm{BKM}}(\Phi_N), N \in \text{Gritsenko index}\}$:
\begin{itemize}
\item At $N = 5$: $24 / N = 24/5 \notin \mathbb{Z}$ so the level-$5$
paramodular index condition (A) fails; while the Borcherds lift
$\mathrm{BL}(\phi_5)$ exists, it does not lie in
$M_{2}(\Gamma_5^{\mathrm{para}})$ at index $1$ without polar saturation
\textup{(Cl\'ery-Gritsenko-Hulek 2015 Section~3 explicitly excludes
$N = 5, 7$ from the Gritsenko paramodular census)}.
\item At $N = 8$: the Frame shape $1^2 2^1 4^1 8^2$ has \emph{non-uniform}
divisor-cycle multiplicities ($m_2 = m_4 = 1$, $m_1 = m_8 = 2$);
the orbifold averaging produces $\phi_8$ as a weak Jacobi form for
$\Gamma_0(8)$, but the Borcherds lift output is a Siegel modular form
of weight $1$ on $\Gamma_8 \subset \mathrm{Sp}_4(\mathbb{Q})$ that does
not match the Gritsenko paramodular family of $\Delta_w$ forms (which
require uniform divisor-cycle multiplicities matching the
``$1$-cycles + $N$-cycles only'' pattern at $N = 2, 3, 6$ or the
balanced pattern $1^4 2^2 4^4$ at $N = 4$).
\end{itemize}
The Gritsenko-paramodular indices are therefore exactly
$\{1, 2, 3, 4, 6\}$ \textup{(Gritsenko 1994 Theorem~3 enumerating the
Borcherds-paramodular weights $\{5, 4, 3, 2, 1\}$ on $\Gamma_N$ for
$N \in \{1, 2, 3, 4, 6\}$; Cl\'ery-Gritsenko 2013 Theorem~1.1 extending
to $N = 4$; Cl\'ery-Gritsenko-Hulek 2015 Theorem~1.2 closing the
$N = 6$ case)}.

\textit{The arithmetic explanation.}  The Gritsenko paramodular indices
$\{1, 2, 3, 4, 6\}$ are exactly the divisors of
$12 = \mathrm{lcm}\{1, 2, 3, 4, 6\}$ that lift to symplectic K3
automorphisms with isolated fixed points (Mukai 1988) and that have
uniform divisor-cycle multiplicities in the Frame shape (i.e., Frame
shape supported only on cycle-lengths that are divisors of $N$, with
$m_a = 24/(N \cdot d(N))$ for $a \in \mathrm{Div}(N) \setminus \{1\}$
where $d(N)$ is the number of nontrivial divisors of $N$).
This intersection of three arithmetic conditions (divisibility of $12$
by $N$, Mukai classification, uniform Frame-shape multiplicities)
selects exactly the five Gritsenko paramodular indices.
\end{proof}

\begin{remark}[Falsification provenance and primary-source attributions]
\label{rem:gritsenko-regularity-attribution}
The proposition combines three independent classification results:
\textsc{Mukai 1988} \textup{(symplectic K3 automorphism classification,
$N \le 8$, eleven distinct conjugacy classes in $M_{24}$);}
\textsc{Gritsenko 1994 + Gritsenko-Nikulin 1998 + Cl\'ery-Gritsenko 2013
+ Cl\'ery-Gritsenko-Hulek 2015} \textup{(paramodular Borcherds-lift
construction at indices $\{1, 2, 3, 4, 6\}$ with weights
$\{5, 4, 3, 2, 1\}$);} and \textsc{Gaberdiel-Hohenegger-Volpato 2010 +
Eguchi-Hikami 2011} \textup{(twined elliptic-genus computation of
$c_N(0)$ for the eight $M_{24}$ classes $N\mathrm{A}/N\mathrm{B}$).}

The original contribution of this proposition is to identify the three
arithmetic conditions (A), (B), (C) of which the intersection is
exactly the Gritsenko-paramodular index list, and to ascribe the
$N = 5, 7, 8$ exclusions to specific Frame-shape and Atkin-Lehner
mechanisms.  Each of (A), (B), (C) is independently checkable; the
selection $\{1, 2, 3, 4, 6\}$ as their intersection is the
\textbf{lattice-regularity criterion}.

The complementary identity at $N = 5, 8$ ($c_N(0) = \chi(K3^{g_N})$
holds numerically but the lift exits the Gritsenko family) reveals
that the Mukai-Lefschetz halving condition (item~(C) above) is
\emph{necessary but not sufficient} for the Gritsenko-paramodular
membership, supplying a partial explanation for why the K3-lattice
classification of symplectic automorphisms is finer than the
Borcherds-paramodular census.
\end{remark}

\subsection{Three constructions, three chiral algebras from K3}
\label{subsec:k3-three-algebraizations}  % label preserved for cross-references

% FOUR OBJECTS:
% 1. B(A) = bar coalgebra = T^c(s^{-1} A-bar) with deconcatenation + twist
% 2. A^i = H^*(B(A)) = dual coalgebra (Koszul cohomology of bar)
% 3. A^! = ((A^i)^v) = dual ALGEBRA (linear or Verdier dual)
% 4. Z^der_ch(A) = derived chiral center = Hochschild cochains = bulk

The failure of $\chi_{\mathrm{top}}/24$ as a modular characteristic is not an isolated
numerical accident. It is the first symptom of a trichotomy that lies at the heart
of the CY programme: three \emph{different mathematical constructions} each extract
a chiral algebra from K3 geometry, and the modular
characteristic $\kappa_{\mathrm{ch}}$ distinguishes them completely.
Only one of these constructions is the CY-to-chiral correspondence $\Phi_2$; the other two
are independent constructions (orbifold, lattice VOA) that do not factor
through~$\Phi_2$.

\begin{proposition}[The K3 trichotomy]
\label{prop:k3-trichotomy}
%: all kappa occurrences carry subscripts (kappa_ch, kappa_cat, kappa_BKM, kappa_fiber)
Let $S$ be a K3 surface with $\chi_{\mathrm{top}}(S) = 24$ and $\chi(\cO_S) = 2$. Three independent constructions produce three chiral algebras with distinct modular characteristics:
\begin{center}
\renewcommand{\arraystretch}{1.4}
\begin{tabular}{llccl}
 \toprule
 \textbf{Route} & \textbf{Algebra} & $c$ & $\kappa_{\mathrm{ch}}$ & \textbf{Source of $\kappa_{\mathrm{ch}}$} \\
 \midrule
 CY correspondence $\Phi_2$ & $\Phi_2(D^b(\Coh(S)))$ & $24$ & $2$ & $\chi(\cO_S) = h^{0,0} - h^{0,1} + h^{0,2}$ \\
 Monster orbifold & $V^\natural$ & $24$ & $12$ & $\Z/2$-orbifold of rank-$24$ lattice VOA \\
 Lattice VOA & $V_\Lambda$ (Leech) & $24$ & $24$ & $\operatorname{rank}(\Lambda) = 24$ \\
 \bottomrule
\end{tabular}
\end{center}
All three share $c = 24$ and originate from K3 geometry (the connection of $V^\natural$ to K3 is indirect: via the Leech lattice and the Niemeier classification). The three modular characteristics $2$, $12$, $24$ are pairwise distinct; no two agree.
\ClaimStatusProvedHere
\end{proposition}

\begin{proof}
(i) \emph{CY functor.} The functor
$\Phi \colon D^b(\Coh(S)) \to \mathrm{ChirAlg}$ of Theorem~\ref{thm:cy-to-chiral}
produces a chiral algebra with $24$ generators (one per Hochschild class, since
$\dim \HH_*(D^b(S)) = 24$) and central charge $c = 24$.
The modular characteristic is $\kappa_{\mathrm{ch}}(\Phi(D^b(S))) = \chi(\cO_S) = 2$:
the CY trace on $\HH_0$ picks out the holomorphic Euler characteristic, not the
total Hochschild dimension. Five independent verifications appear in
Theorem~\ref{thm:k3-kappa}.

(ii) \emph{Monster orbifold.} The moonshine module $V^\natural$ has $c = 24$,
constructed as the $\Z/2$-orbifold of the Leech lattice VOA $V_\Lambda$.
Its modular characteristic is $\kappa_{\mathrm{ch}}(V^\natural) = 12$
(Vol~III, Remark~\ref{rem:lattice-voa-k3} and the lattice VOA bar complex of
Vol~I). The connection to K3 is that the Leech lattice itself is the unique
even unimodular lattice of rank $24$ with no roots, and its root system is
controlled by the K3 geometry: the $196560$ minimal vectors of $\Lambda_{\mathrm{Leech}}$
arise as a shadow of $\chi_{\mathrm{top}}(S) = 24$. The $\Z/2$-orbifold halves
$\kappa_{\mathrm{ch}}$ from $24$ to $12$.

(iii) \emph{Lattice VOA.} The Leech lattice VOA $V_\Lambda$ has $c = 24$ and
$\kappa_{\mathrm{ch}}(V_\Lambda) = \operatorname{rank}(\Lambda) = 24$
(Remark~\ref{rem:leech-connection}). As a lattice VOA, it is the free-field
realization of $24$ bosons compactified on $\Lambda_{\mathrm{Leech}}$. The K3
lattice $\Lambda_{K3} = U^3 \oplus E_8(-1)^2$ of rank $22$ and signature $(3,19)$
embeds in the Leech lattice after adding a hyperbolic plane $U$; the passage
$\Lambda_{K3} \oplus U \hookrightarrow \Lambda_{\mathrm{Leech}}$ is the Niemeier
classification step that ties K3 geometry to the rank-$24$ lattice.
\end{proof}

The three numbers $2$, $12$, $24$ are not random.
They encode the depth at which each construction interrogates the K3 geometry:
\begin{itemize}
\item $\kappa_{\mathrm{ch}} = 2 = \chi(\cO_S)$: the CY functor sees only the
 holomorphic Euler characteristic, the coarsest algebraic invariant of the surface.
 The $24$ Hochschild generators are present, but the CY trace projects onto a
 two-dimensional subspace.
\item $\kappa_{\mathrm{ch}} = 12$: the Monster orbifold $V^\natural$ is the
 $\Z/2$-orbifold of $V_\Lambda$. The orbifolding halves $\kappa_{\mathrm{ch}}$
 from $24$ to $12$, as the $\Z/2$-twisted sector kills exactly half of the
 generators that contribute to the shadow tower.
\item $\kappa_{\mathrm{ch}} = 24 = \dim \HH_*(D^b(S))$: the lattice VOA sees
 the full Hochschild homology. Every class contributes a free boson, and
 $\kappa_{\mathrm{ch}} = \operatorname{rank}$ for lattice algebras
 (Vol~I, census entry C1 applied at rank~$24$).
\end{itemize}
The ratio $\kappa_{\mathrm{ch}}(V_\Lambda) / \kappa_{\mathrm{ch}}(V^\natural) /
\kappa_{\mathrm{ch}}(\Phi(D^b(S))) = 24 : 12 : 2 = 12 : 6 : 1$
is an invariant of the construction, not of K3 itself. The K3 surface
supplies a single derived category $D^b(\Coh(S))$ and a single Hochschild homology
$\HH_*(D^b(S))$ of total dimension~$24$; the three constructions differ in how much of
this $24$-dimensional space they promote to a chiral algebra generator with
nonvanishing contribution to the shadow tower.

\begin{remark}[Why one geometry produces three modular characteristics]
\label{rem:k3-trichotomy-why}
%: kappa_ch subscripted throughout
The trichotomy is not a failure of the CY-to-chiral correspondence programme. It is evidence that
the correspondence $\Phi_2$ is \emph{not} the only route from geometry to chiral algebra,
and that $\kappa_{\mathrm{ch}}$ detects the route, not merely the destination.

At the level of higher structures, the three constructions differ in the operadic
data they extract. The CY correspondence $\Phi_2$ uses the $\Eone$-algebra structure on the
bar complex $B^{\Eone}(D^b(S))$, which retains the ordering of collision points;
the projection to the shadow tower (the $\Sigma_n$-coinvariant quotient) loses all
but the trace, producing $\kappa_{\mathrm{ch}} = 2$. The lattice VOA
construction uses the $E_\infty$-structure directly, promoting every lattice vector
to a generator; no information is lost in the passage from Hochschild to chiral,
and $\kappa_{\mathrm{ch}} = 24$. The Monster orbifold sits between:
the $\Z/2$-orbifolding is a structured projection that preserves exactly half the
data, giving $\kappa_{\mathrm{ch}} = 12$.

The question this raises for the CY programme is sharp: given a CY$_d$
variety $X$, is there a \emph{canonical} construction, or do multiple
independent constructions produce distinct chiral algebras from the
same geometry? For $d = 1$ (elliptic curve) the question is trivial:
$\kappa_{\mathrm{ch}}^{\mathrm{Heis}}(E) = 1$ by rank-additivity on all Heisenberg-level presentations. For $d = 2$ (K3) the question is
the trichotomy of this subsection. For $d = 3$ ($K3 \times E$) the question
acquires a fourth layer: $\kappa_{\mathrm{BKM}} = 5$, the Borcherds lift weight,
is not the modular characteristic of any of the three $d = 2$ algebras but of the
second-quantized BKM superalgebra $\mathfrak{g}_{\Delta_5}$ (Proposition~\ref{prop:k3e-fiber-global}).
\end{remark}

\noindent\textit{Verification}: 75 tests in \texttt{k3\_cy\_a2\_verification\_engine.py} covering the three-construction distinction, $\kappa_{\mathrm{ch}}$ values, orbifold halving, and $\dim \HH_*$-to-$\kappa_{\mathrm{ch}}$ comparison for all three constructions.



\section{Six routes to $G(K3 \times E)$}
\label{sec:k3e-six-routes}

The quantum vertex chiral group $G(K3 \times E)$ is approached by six independent mathematical constructions. Only Route~4 uses the CY-to-chiral correspondence $\Phi_3$; the remaining five are independent constructions. That all six converge on the same target is the \emph{content} of Conjecture~CY-C, not a consequence of functoriality.

\begin{enumerate}[label=\textbf{Route \arabic*.}, leftmargin=*]
 \item \textbf{BLLPRR (Bryan--Leung--Lian--Pandharipande--Ruan--R).} Start from the DT theory of $K3 \times E$ (Theorem~\ref{thm:dt-igusa}). The DT partition function $Z^X = C / \Delta_5^2$ gives $\mathfrak{g}_{\Delta_5}$ via the product formula (Theorem~\ref{thm:k3e-product}). The CoHA (Theorem~\ref{thm:k3e-coha}) gives $U(\fn_+)$. The braiding comes from the MO $R$-matrix (Theorem~\ref{thm:k3e-mo-rmatrix}). \emph{Status}: proved (modulo standard conjectures on DT/GW for $K3 \times E$).

 \item \textbf{Holomorphic-topological at generic K3.} At a generic point of $\cM_{K3}$, the K3 sigma model has no enhanced symmetry and the $20$ left-moving bosons decouple into $20$ Heisenberg factors $\cH_1^{\oplus 20}$. The $K3 \times E$ theory is then $(20$ Heisenberg$) \otimes (\text{elliptic theory})$, giving the positive half as $\CoHA \simeq U(\hat{\mathfrak{h}}_{20}^+)$ (abelian, class G). The full $\mathfrak{g}_{\Delta_5}$ appears only after automorphic completion. \emph{Status}: proved at the level of free-field realization; automorphic completion is Route 1.

 \item \textbf{Holomorphic-topological at enhanced K3.} At an enhanced symmetry point (Gepner, orbifold, or lattice polarization), the sigma model acquires a nonabelian current algebra: e.g.\ at the $T^4/\Z_2$ orbifold point, one gets $V_1(\mathfrak{so}_4)^{\oplus 4}$ (level-$1$ affine $\mathfrak{so}_4$). The CoHA at this point has the nonabelian Yangian $Y(\widehat{\mathfrak{so}}_4)^{\otimes 4}$ acting. The passage from enhanced to generic is a deformation (wall-crossing in $\cM_{K3}$). \emph{Status}: proved at specific orbifold points; deformation to generic is Route 2.

 \item \textbf{CY$_3$ quantum vertex chiral group (Theorem CY-A$_3$).} Apply the abstract CY-to-chiral correspondence $\Phi_3 \colon D^b(\Coh(K3 \times E)) \to \text{ChirAlg}$ of Theorem~CY-A$_3$ (Theorem~\ref{thm:cy-to-chiral-d3}). This gives $A_{K3 \times E} = \Phi_3(D^b(\Coh(K3 \times E)))$ directly. The quantum vertex chiral group $G(K3 \times E) = G(A_{K3 \times E})$ requires Conjecture~CY-C.  \emph{Status}: the chiral algebra $A_{K3 \times E}$ is constructed (CY-A$_3$ proved); the group object $G(K3 \times E)$ is conjectural (CY-C).

 \item \textbf{AdS$_3$/CFT$_2$.} The type IIB string on $\mathrm{AdS}_3 \times S^3 \times K3$ at $k$ units of flux has boundary CFT$_2$ containing the symmetric orbifold $\Sym^k(K3)$ in the large-$k$ limit. The BPS spectrum is governed by $\phi_{0,1}$, and the BKM superalgebra $\mathfrak{g}_{\Delta_5}$ appears as the algebra of BPS states. The AdS$_3$ Yangian is the bulk-boundary manifestation of the MO $R$-matrix. \emph{Status}: heuristic (the bulk-boundary dictionary is not mathematically rigorous, but the BPS counting is exact by supersymmetry).

 \item \textbf{BLLPR Schur sector.} Start from the $6$d $(2,0)$ theory of type~$A_1$ on $K3$. Compactify on a Riemann surface~$C$ to obtain a $4$d $\cN = 2$ class~$S$ theory $T[A_1, C]$. The BLLPR correspondence (Beem--Lemos--Liendo--Peelaers--Rastelli, arXiv:1312.5344) associates to this $4$d theory a $2$d chiral algebra on~$C$: the $\cW$-algebra $\cW_k(\fsl_2)$ at a level~$k$ determined by the pants decomposition of~$C$.  For $\fsl_2$ (rank~$1$), $\cW_k(\fsl_2) = \mathrm{Vir}_{c}$ with $c = 13 - 6(k+2) - 6/(k+2)$.  \emph{Status}: the BLLPR correspondence is proved (arXiv:1312.5344); the connection to the K3 Yangian is conjectural.
\end{enumerate}

\begin{remark}[Route comparison]
\label{rem:k3e-route-comparison}
Routes 1--3 give progressively more detailed structural information: Route 1 captures the combinatorial skeleton ($\fn_+$ and the root system), Route 2 gives the generic (abelian) fiber, and Route 3 gives the enhanced (nonabelian) fiber. Route 4 gives the full chiral algebra $A_{K3 \times E}$ via Theorem~CY-A$_3$. Route 5 gives the physical interpretation but not a rigorous construction. Route~6 gives the $4$d $\cN = 2$ Schur-sector perspective. The relative chiral algebra of Section~\ref{sec:k3e-relative-chiral} provides a concrete intermediary between Routes 2--3 and Route 4.
\end{remark}

\begin{remark}[The BLLPR algebra and the K3 Yangian are distinct algebraizations]
\label{rem:bllpr-k3-connection}
Route~6 produces a chiral algebra ($\cW_k(\fsl_2)$) on the \emph{same} curve~$C$ as the K3 Yangian of Section~\ref{subsec:k3-yangian}.  These are \textbf{different algebraizations} of the same underlying $6$d physics, distinguished by five structural invariants:
\begin{enumerate}[label=\textup{(\roman*)}]
 \item \emph{Central charge.} $\cW_k(\fsl_2)$ has $c = 13 - 6(k+2) - 6/(k+2)$ (e.g.\ $c = -2$ at $k = -3/2$). The K3 Heisenberg has $c = 24$ (Mukai rank).  The Sugawara construction inside the rank-$24$ Heisenberg produces a Virasoro subalgebra at $c_{\mathrm{Sug}} = 24\Psi/(\Psi + 1)$; setting $c_{\mathrm{Sug}} = c_{\mathrm{BLLPR}}$ determines the effective level $\Psi = c_{\mathrm{BLLPR}}/(24 - c_{\mathrm{BLLPR}})$.  At $k = -3/2$: $\Psi = -1/13$.

 \item \emph{$E_n$ structure.} The BLLPR algebra $\cW_k(\fsl_2)$ is a vertex algebra ($E_\infty$-chiral). The K3 Yangian is $\Eone$-chiral (ordered: the $E_\infty$ averaging map kills the Hopf structure).

 \item \emph{$\Omega$-dependence.} The BLLPR chiral algebra does \emph{not} depend on the $\Omega$-background: it is constructed from the Schur sector of the physical $4$d theory. The K3 Yangian \emph{requires} the $\Omega$-background (the deformation parameter $\sigma_3 = h_1 h_2 h_3$; at $\sigma_3 = 0$ the Yangian degenerates to the undeformed Heisenberg). The $\Omega$-background is the intermediary between the two.

 \item \emph{Character growth.}  The BLLPR vacuum character at energy~$n$ counts partitions of~$n$ into parts $\geq 2$ (dimension $p_{\geq 2}(n)$: $1, 0, 1, 1, 2, 2, 4, \ldots$).  The K3 Fock character counts $24$-coloured partitions ($p_{24}(n)$: $1, 24, 324, 3200, \ldots$).  The ratio $p_{24}(n)/p_{\geq 2}(n)$ grows exponentially.

 \item \emph{Modular group.}  BLLPR characters transform under $\SL_2(\bZ)$; the K3 Yangian characters involve $\Sp_4(\bZ)$ (Siegel modular forms, Section~\ref{sec:k3e-borcherds-lift}).
\end{enumerate}

\noindent\textbf{Conjectural connection.}  The BLLPR $\cW$-algebra is the $\Omega$-independent sector of the K3 Yangian: passing to the cohomology of the Schur supercharge $Q_{\mathrm{Schur}}$ projects the full K3 Yangian Fock space onto the BLLPR vacuum module.  The Sugawara embedding at $\Psi = c_{\mathrm{BLLPR}}/(24 - c_{\mathrm{BLLPR}})$ realizes this projection.  The chiral algebra $A_{K3 \times E}$ exists by Theorem~CY-A$_3$; the K3 Yangian as a full quantum group object requires Conjecture~CY-C.
\end{remark}

\noindent\textit{Verification}: 73 tests in \texttt{bllpr\_k3\_connection.py} covering the BLLPR central charge formula at $5$ levels, level-rank duality, Virasoro vacuum character, K3 Fock character ($1/\eta^{24}$) through $q^6$, Sugawara embedding roundtrip, $\Omega$-deformation analysis, character growth comparison, compactification chain ($6$d $\to$ $4$d $\to$ $2$d), Route~$6$ structural distinction, and $4$ cross-engine checks.

%% ========================================================================

\section{Enhanced symmetry points of K3 moduli and the Niemeier landscape}
\label{sec:k3-moduli-niemeier-landscape}

At the generic point of the K3 moduli space
$\cM_{K3} = \mathrm{O}(3,19;\bZ) \backslash \mathrm{O}(3,19;\bR) / (\mathrm{O}(3) \times \mathrm{O}(19))$,
the chiral algebra of the K3 sigma model is the rank-$24$ Mukai
Heisenberg $\cH_{\mathrm{Muk}}$ with pairing of signature $(4,20)$,
central charge $c = 24$, and modular characteristic
$\kappa_{\mathrm{ch}} = 2 = \chi(\cO_{K3})$.
At special loci (the \emph{enhanced symmetry points}) the
chiral algebra acquires nonabelian current subalgebras, and the
representation theory becomes richer.  This section classifies these
loci via the Niemeier lattice landscape and verifies that
$\kappa_{\mathrm{ch}}$ is invariant across the entire moduli space.

\subsection{ADE enhancement at singular points}
\label{subsec:k3-ade-enhancement}

Near an isolated ADE singularity of type $\fg$ on K3, the chiral
algebra of the sigma model enhances to include an affine
Kac--Moody subalgebra $\hat{\fg}_1$ at level~$1$.

\begin{proposition}[ADE enhancement classification]
\label{prop:k3-ade-enhancement}
\ClaimStatusProvedElsewhere
Let $S$ be a K3 surface acquiring an isolated ADE singularity of
type $\fg$, with $\fg$ a simply-laced simple Lie algebra of rank~$r$
and dual Coxeter number~$h^\vee$.
The K3 sigma model at this point contains the affine subalgebra
$\hat{\fg}_1$ with:
\begin{center}
\renewcommand{\arraystretch}{1.3}
\begin{tabular}{lcccc}
\toprule
\textbf{Singularity} & \textbf{Subalgebra} & $\rank$ & $c(\hat{\fg}_1)$ & \textbf{Shadow class} \\
\midrule
$A_n$ & $\widehat{\fsl}_{n+1,1}$ & $n$ & $n$ & $L$ \\
$D_n$ & $\widehat{\fso}_{2n,1}$ & $n$ & $n$ & $L$ \\
$E_6$ & $\hat{E}_{6,1}$ & $6$ & $6$ & $L$ \\
$E_7$ & $\hat{E}_{7,1}$ & $7$ & $7$ & $L$ \\
$E_8$ & $\hat{E}_{8,1}$ & $8$ & $8$ & $L$ \\
\bottomrule
\end{tabular}
\end{center}
At level~$1$, the central charge of $\hat{\fg}_1$ equals the rank:
$c(\hat{\fg}_1) = \dim(\fg)/(1 + h^\vee) = \rank(\fg)$ for all
simply-laced~$\fg$.  The shadow class of the enhanced subalgebra is~$L$
\textup{(}linear, shadow depth~$3$\textup{)}: the cubic shadow $S_3 \neq 0$
\textup{(}from the structure constants\textup{)} but the quartic shadow
$S_4 = 0$ at level~$1$, giving critical discriminant $\Delta = 0$.
\end{proposition}

\begin{remark}[Shadow class $L$ for the subalgebra, class $M$ for the full sigma model]
\label{rem:k3-ade-shadow-subalgebra-vs-full}
The shadow class of the \emph{enhanced subalgebra} $\hat{\fg}_1$ is~$L$
(depth~$3$), but the shadow class of the \emph{full K3 sigma model}
at an ADE point remains~$M$ (infinite depth, as in
Computation~\ref{comp:shadow-tower-k3}).  The enhancement adds
a nonabelian current sector with $S_3 \neq 0$, but the infinite
shadow depth of the full sigma model arises from the
$\cN = 4$ superconformal structure interacting with the
Hochschild data, not from the enhanced subalgebra alone.
The shadow class computation for the full algebra follows
from the shadow tower data $(\kappa_{\mathrm{ch}}, S_3, S_4) = (2, 2, 5/156)$
of Computation~\ref{comp:shadow-tower-k3}, which is
\emph{moduli-independent}: the $(S_3, S_4)$ values are determined
by the $\cN = 4$ Ward identities, not by the specific enhancement.
\end{remark}

\subsection{Moduli invariance of $\kappa_{\mathrm{ch}}$}
\label{subsec:k3-kappa-moduli-invariance}

\begin{theorem}[$\kappa_{\mathrm{ch}}$ is constant across K3 moduli]
\label{thm:k3-kappa-moduli-invariance}
\ClaimStatusProvedHere
The modular characteristic of the K3 sigma model is
\[
 \kappa_{\mathrm{ch}}(\cA_{K3}) = 2 = \chi(\cO_{K3})
\]
at every point of $\cM_{K3}$: generic, ADE-enhanced, Kummer,
Gepner, and Niemeier.
\end{theorem}

\begin{proof}
The modular characteristic $\kappa_{\mathrm{ch}}$ depends on the
chiral algebra structure and is a deformation invariant:
it is computed from the genus-$1$ obstruction of the bar complex,
which depends only on the $\cN = 4$ superconformal structure, not
on the specific point in moduli.  Five paths verify this:

(i) \emph{Geometric.} $\kappa_{\mathrm{ch}} = \chi(\cO_S) =
h^{0,0} - h^{0,1} + h^{0,2} = 1 - 0 + 1 = 2$ is a topological
invariant of K3, computed from the Hodge numbers which are
constant on the moduli space.

(ii) \emph{Kummer.} At $K3 = T^4/\bZ_2$:
$\kappa_{\mathrm{ch}}(\cA^{\mathrm{orb}}) = 2$
(Proposition~\ref{prop:kummer-orbifold}).  The $8$ untwisted
Heisenberg generators contribute $0$; the $16$ twisted sectors
produce the level-$1$ affine $\mathfrak{so}_4^{\oplus 4}$
enhancement with $\kappa_{\mathrm{ch}} = 4 \times \frac{1}{2} = 2$.

(iii) \emph{ADE invariance.} At an ADE singularity of type~$\fg$,
the enhanced subalgebra $\hat{\fg}_1$ has its own
$\kappa_{\mathrm{ch}}(\hat{\fg}_1) = \rank(\fg)$, but the full sigma model
retains $\kappa_{\mathrm{ch}} = 2$ because the $\cN = 4$ Ward identities
constrain the genus-$1$ obstruction.  The enhanced subalgebra replaces
part of the rank-$24$ Heisenberg (which contributes $\kappa_{\mathrm{ch}} = 24$
as a standalone lattice VOA) with a nonabelian algebra at the same
central charge; the reduction from $24$ to $2$ is due to the
$\cN = 4$ superconformal projection, which is present at all moduli.

(iv) \emph{Gepner.} At the Fermat quartic point, the Gepner tensor
product of four $\cN = 2$ minimal models gives
$\kappa_{\mathrm{Gepner}} = 3$ (Remark~\ref{rem:gepner-model-k3}), but this
is a \emph{different algebraization route}: the physical sigma model
at the Gepner point has the full $\cN = 4$ structure with
$\kappa_{\mathrm{ch}} = 2$, not~$3$.

(v) \emph{Wall-crossing.}  Moving between moduli points crosses walls
in the Bridgeland stability manifold (Section~\ref{sec:k3e-wall-crossing}).  These
wall-crossings correspond to flops, which preserve $\kappa_{\mathrm{ch}}$ (
flops are autoequivalences, $\kappa_{\mathrm{ch}}(A_X) = \kappa_{\mathrm{ch}}(A_{X^+})$).  In
particular, $\kappa_{\mathrm{ch}}$ cannot jump between enhancement points.
\end{proof}

\noindent\textit{Verification}: $20$ tests in
\texttt{niemeier\_shadow\_landscape.py} (\texttt{TestKappaCHModuliIndependence}
and \texttt{TestMultiPathCrossChecks::test\_kappa\_ch\_k3\_three\_paths}) covering
$\kappa_{\mathrm{ch}} = 2$ at generic, Kummer, Gepner, and all ADE singularity points
$A_1$--$A_8$, $D_4$--$D_8$, $E_6$, $E_7$, $E_8$; plus three-path cross-check
via $\chi(\cO_{K3})$, $\dim_{\bC}(K3)$, and Kummer orbifold.

\subsection{The 24 Niemeier lattices and maximal enhancement}
\label{subsec:k3-niemeier-landscape}

The maximal enhancement points of K3 moduli are classified by
the $24$ Niemeier lattices: the even unimodular lattices of rank~$24$.
At such a point, the full Mukai lattice
$\widetilde{\Lambda}_{K3} = U^4 \oplus E_8(-1)^2$ of
Remark~\ref{rem:lattice-voa-k3} admits a root system $R$ from one
of the $24$ Niemeier root systems, and the lattice VOA $V_N$
of the Niemeier lattice $N$ provides the maximal enhanced algebra.

\begin{proposition}[Niemeier shadow landscape]
\label{prop:niemeier-shadow-landscape}
\ClaimStatusProvedHere
Let $N$ be one of the $24$ Niemeier lattices with root system $R(N)$.
The lattice VOA $V_N$ has:
\begin{center}
\renewcommand{\arraystretch}{1.3}
\begin{tabular}{ccccl}
\toprule
$c(V_N)$ & $\kappa_{\mathrm{ch}}(V_N)$ & \textbf{Class} & \textbf{Depth} & \textbf{Criterion} \\
\midrule
$24$ & $24$ & $G$ & $2$ & $S_3 = S_4 = 0 \implies \Delta = 0$ \\
\bottomrule
\end{tabular}
\end{center}
These values are \emph{identical} for all $24$ Niemeier lattices.
The shadow obstruction tower is blind to the root system~$R(N)$:
it cannot distinguish the Leech lattice \textup{(}$R = \varnothing$\textup{)}
from $A_1^{24}$ \textup{(}$|R| = 48$\textup{)} or $E_8^3$ \textup{(}$|R| = 720$\textup{)}.
\end{proposition}

\begin{proof}
Lattice VOAs are $E_\infty$-algebras (commutative), so the cubic shadow
$S_3 = 0$ (no structure constants contribute to the bar complex at
degree~$3$).  The quartic shadow $S_4 = 0$ since the lattice VOA is a
free-field realization with no Sugawara-type corrections.
The critical discriminant $\Delta = 8\kappa_{\mathrm{ch}} S_4 = 0$, so the
single-line dichotomy theorem (Vol~I,
Theorem~\ref{thm:riccati-algebraicity}) places all lattice VOAs in
class~$G$ with shadow depth~$2$.  The shadow class depends only on
$(\kappa_{\mathrm{ch}}, S_3, S_4) = (24, 0, 0)$; the root system $R(N)$
appears neither in $\kappa_{\mathrm{ch}}$ (which is the lattice rank) nor
in $S_3$ or $S_4$ (which vanish for any lattice VOA).  \emph{This is
 compliant}: the shadow class is determined from the full tower
computation $S_3 = S_4 = 0 \Rightarrow \Delta = 0 \Rightarrow$ class~$G$,
not from generator counting.
\end{proof}

\begin{remark}[Distinguishing Niemeier lattices: the theta series]
\label{rem:niemeier-theta-distinguishing}
Since the shadow tower cannot distinguish the $24$ Niemeier lattices,
the distinguishing invariant is the theta series.  Every Niemeier lattice
$N$ has
\[
 \Theta_N(\tau) = E_{12}(\tau) + c_\Delta(N) \cdot \Delta(\tau),
\]
where $E_{12}$ is the weight-$12$ Eisenstein series and $\Delta$ is the
Ramanujan delta function.  The coefficient $c_\Delta$ is determined by
the number of roots $|R(N)|$:
\[
 c_\Delta(N) = \frac{691 \cdot |R(N)| - 65520}{691}.
\]
\end{remark}

The following table records $|R(N)|$, $c_\Delta$, and the $q^2$ theta
coefficient (the number of vectors of norm~$4$) for all~$24$
Niemeier lattices.  In every case the lattice VOA shadow data is
$(\kappa_{\mathrm{ch}}, \text{class}, \text{depth}) = (24, G, 2)$ and the K3
sigma model has $\kappa_{\mathrm{ch}} = 2$.

\begin{center}
\small
\renewcommand{\arraystretch}{1.2}
\begin{tabular}{rlrrr}
\toprule
\# & \textbf{Root system} & $|R|$ & $c_\Delta$ & $\Theta_N|_{q^2}$ \\
\midrule
$1$ & Leech (no roots) & $0$ & $-65520/691$ & $196560$ \\
$2$ & $D_{24}$ & $1104$ & $697344/691$ & $170064$ \\
$3$ & $D_{16}\,E_8$ & $720$ & $432000/691$ & $179280$ \\
$4$ & $E_8^3$ & $720$ & $432000/691$ & $179280$ \\
$5$ & $A_{24}$ & $600$ & $349080/691$ & $182160$ \\
$6$ & $D_{12}^2$ & $528$ & $299328/691$ & $183888$ \\
$7$ & $A_{17}\,E_7$ & $432$ & $232992/691$ & $186192$ \\
$8$ & $D_{10}\,E_7^2$ & $432$ & $232992/691$ & $186192$ \\
$9$ & $A_{15}\,D_9$ & $384$ & $199824/691$ & $187344$ \\
$10$ & $D_8^3$ & $336$ & $166656/691$ & $188496$ \\
$11$ & $A_{12}^2$ & $312$ & $150072/691$ & $189072$ \\
$12$ & $A_{11}\,D_7\,E_6$ & $288$ & $133488/691$ & $189648$ \\
$13$ & $E_6^4$ & $288$ & $133488/691$ & $189648$ \\
$14$ & $A_9^2\,D_6$ & $240$ & $100320/691$ & $190800$ \\
$15$ & $D_6^4$ & $240$ & $100320/691$ & $190800$ \\
$16$ & $A_8^3$ & $216$ & $83736/691$ & $191376$ \\
$17$ & $A_7^2\,D_5^2$ & $192$ & $67152/691$ & $191952$ \\
$18$ & $A_6^4$ & $168$ & $50568/691$ & $192528$ \\
$19$ & $D_4^6$ & $144$ & $33984/691$ & $193104$ \\
$20$ & $A_5^4\,D_4$ & $144$ & $33984/691$ & $193104$ \\
$21$ & $A_4^6$ & $120$ & $17400/691$ & $193680$ \\
$22$ & $A_3^8$ & $96$ & $816/691$ & $194256$ \\
$23$ & $A_2^{12}$ & $72$ & $-15768/691$ & $194832$ \\
$24$ & $A_1^{24}$ & $48$ & $-32352/691$ & $195408$ \\
\bottomrule
\end{tabular}
\end{center}

\noindent\textit{Verification}: $72$ tests in
\texttt{niemeier\_shadow\_landscape.py} covering all $24$
Niemeier lattices (rank, root count, shadow data, theta
coefficients, $c_\Delta$ computation), $\kappa_{\mathrm{ch}}$ moduli
independence at $19$ enhancement points, and $12$ multi-path
cross-checks (root count via component formula vs direct formula,
theta decomposition $E_{12} + c_\Delta \cdot \Delta$ self-consistency,
dimension formula cross-checks, level-$1$ central charge from two
independent formulas, $\kappa_{\mathrm{ch}} = 2$ via three independent
paths, and Leech $q^2 = 196560$ via known value vs decomposition
formula).

\begin{conjecture}[Mathieu $M_{24}$ moonshine through the K3 Yangian]
\label{conj:mathieu-moonshine-yangian}
\ClaimStatusConjectured
The Mathieu group $M_{24}$ acts on the K3 Yangian $Y(\frakg_{K3})$ by
permutation of the $24$ parameters $h_1, \ldots, h_{24}$.
\begin{enumerate}[label=\textup{(\alph*)}]
  \item \emph{Parameter action.}
  $M_{24} \hookrightarrow S_{24}$ acts on the parameter space by
  $\sigma \cdot (h_1, \ldots, h_{24}) = (h_{\sigma^{-1}(1)}, \ldots,
  h_{\sigma^{-1}(24)})$, preserving both the CY$_2$ constraint
  $\sum h_i = 0$ and the structure function
  $g_{K3}(z) = \prod (z - h_i)/(z + h_i)$ (which is symmetric in the $h_i$).

  \item \emph{Moonshine factorisation.}
  The Eguchi--Ooguri--Tachikawa coefficients $A_n$ factorise as
  $A_n = \kappa_{\mathrm{ch}} \cdot \dim(\rho_n)$ where
  $\kappa_{\mathrm{ch}} = 2$ and $\rho_n$ are $M_{24}$ representations
  of dimensions $a_n = 45, 231, 770, 2277, \ldots$.
  The mock modular form
  $h(\tau) = 2 q^{-1/8}(-1 + 45q + 231q^2 + 770q^3 + \cdots)$
  has polar coefficient $-\kappa_{\mathrm{ch}} = -2$.

  \item \emph{Twined eta products.}
  For each conjugacy class $[g] \in M_{24}$ with Frame shape
  $\pi_g = \prod_i i^{a_i}$, the twined bar Euler product is the eta
  product $\eta_g(\tau) = \prod_{i} \eta(i\tau)^{a_i}$.
  The identity class~$1A$ gives $\eta(\tau)^{24} = \Delta(q)/q$
  (the Ramanujan discriminant); the free involution~$2B$ gives
  the trivial eta product (trace~$0$).

  \item \emph{Representation decomposition.}
  The charge-$1$ representation $V = \C^{24}$ decomposes as
  $V = V_{23}^{\mathrm{std}} \oplus V_1^{\mathrm{triv}}$
  under $M_{24}$, reflecting the decomposition
  $\Lambda_{\mathrm{Muk}} = \Lambda_{23} \oplus \langle\mathbf{1}\rangle$.
  The twined structure function $g_{[g]}(z)$ has
  degree $(\operatorname{tr}_{V_{24}} g,\; \operatorname{tr}_{V_{24}} g)$.
\end{enumerate}

Of the $25$ conjugacy classes of $M_{24}$, at least~$21$ are realised as
K3 sigma model symmetries (Gaberdiel--Hohenegger--Volpato,
\texttt{arXiv:1008.3778}); the classes $7A$, $7B$, $15A$, $15B$ are not.
The chiral algebraic interpretation links $M_{24}$ to the K3 Yangian
parameter space; it does \emph{not} identify $M_{24}$ as an automorphism
group of $Y(\frakg_{K3})$ (which would require $M_{24}$-equivariance of
the RTT presentation, not merely of the parameters).

The Umbral moonshine programme (Cheng--Duncan--Harvey,
\texttt{arXiv:1204.2779}) extends this to all~$23$ Niemeier root
systems $R(N)$, with $M_{24}$ at the $A_1^{24}$ point (lambency $\ell = 2$)
and distinct Umbral groups $G_N$ at each Niemeier lattice.  The shadow
data at every Niemeier point satisfies
$\kappa_{\mathrm{ch}}^{\text{lattice}} = 24$ and
$\kappa_{\mathrm{ch}}^{\text{K3 }\sigma} = 2$
(Proposition~\ref{prop:niemeier-shadow-landscape}).
\end{conjecture}

\begin{remark}[Classical attribution for the Mathieu frame shape story]
\label{rem:mathieu-moonshine-frame-shape-classical}
The entire Mathieu frame-shape narrative is classical. The observation
that the K3 elliptic genus decomposes as a virtual $M_{24}$-character
was made by Eguchi--Ooguri--Tachikawa \cite{EguchiOoguriTachikawa2010},
with the twining characters computed by Cheng, Gaberdiel--Hohenegger--Volpato
\cite{GaberdielHoheneggerVolpato2010}, and Eguchi--Hikami (2010).
The extension to all 23 Niemeier root systems as Umbral moonshine is
Cheng--Duncan--Harvey \cite{ChengDuncanHarvey2014}. The Frame shape
$\pi_g = \prod_i i^{a_i}$ is a classical character-theoretic invariant
going back to Kondo (1983) / Conway--Norton (1979) for the Monster.
Gannon \cite{Gannon16} surveys the moonshine landscape and settles the
existence of the $M_{24}$-module structure on the massive Mathieu
multiplicities. The programme's contribution is \emph{not} the Frame
shape identification or the twined eta products, but the realisation
of $M_{24}$ as permuting the $24$ parameters of the K3 Yangian ---
an algebraic-parameter action, not an automorphism of $Y(\frakg_{K3})$
(which would require RTT-equivariance). The twined bar Euler products
$\eta_g(\tau) = \prod_i \eta(i\tau)^{a_i}$ are the classical
Conway--Norton / Eguchi--Ooguri--Tachikawa twined eta products;
the bar-cobar recasting rewrites them as twined Euler characteristics
of the rank-$24$ Mukai Heisenberg VOA, nothing more. The moonshine
itself (the existence of an honest $M_{24}$-module) is Gannon;
the bar-cobar language supplies a home but not a new theorem.
\end{remark}

\noindent\textit{Verification}: $50$ tests in
\texttt{mathieu\_moonshine\_yangian.py} covering all $25$ Frame shapes
(sum to $24$, trace consistency, cycle length divisibility), moonshine
coefficients $A_1, \ldots, A_9$ (EOT table), twined eta products (identity,
$2A$, $2B$, $3A$, leading powers), $M_{24}$ action on charge-$1$
representation (CY$_2$ preservation, structure function invariance,
$24 = 23 + 1$ decomposition, $\Sym^2$ and $\mathrm{Alt}^2$ traces),
$23$ Umbral Niemeier root systems, trace-of-power formula for all
classes and prime powers, and $6$ cross-verification tests against
\texttt{k3\_yangian.py}.

\begin{conjecture}[Deeper Mathieu--Yangian connection]
\label{conj:mathieu-yangian-deeper}
\ClaimStatusConjectured\quad
The Mathieu--Yangian connection of
Conjecture~\ref{conj:mathieu-moonshine-yangian} extends along four
axes.  All claims are conditional on the existence of $Y(\frakg_{K3})$
and on CY-A$_2$ \textup{(}PROVED at $d=2$\textup{)}.

\medskip\noindent
\textbf{(I) Generator permutation at the $A_1^{24}$ Niemeier point.}
The $24$ Mukai generators span $H^*(K3, \C) = H^0 \oplus H^2 \oplus H^4$,
with dimensions $1 + 22 + 1 = 24$.  The group $M_{24}$ acts on these
generators by permutation ONLY at the $A_1^{24}$ Niemeier point, where
the $24$ Mukai directions are labelled by the $24$ $A_1$ roots.  At this
point, the Gram matrix of the root system is $2\, I_{24}$ restricted to
$\sum h_i = 0$, and every permutation in~$M_{24}$ preserves the Mukai
pairing.  At a generic K3 point, $M_{24}$ does NOT act on $H^*(K3, \C)$
as a permutation group, because the $24$ generators are not equivalent
under the Mukai pairing.

\medskip\noindent
\textbf{(II) Twined bar Euler product = Frame shape eta product.}
For each $g \in M_{24}$ with Frame shape $\pi_g = \prod_i i^{a_i}$,
the twined bar Euler product
\begin{equation}\label{eq:twined-bar-euler-cyclotomic}
\prod_{n \geq 1} \det\bigl(1 - g \cdot q^n \,\big|\, V_{24}\bigr)
\;=\; \prod_{i \mid a_i \neq 0} \prod_{n \geq 1} (1 - q^{in})^{a_i}
\;=\; \prod_{i} \eta(i\tau)^{a_i}
\;=\; \eta_g(\tau),
\end{equation}
where the first equality is the cyclotomic identity: a $k$-cycle
contributes $\prod_{j=0}^{k-1}(1 - \omega_k^j x) = 1 - x^k$, applied
term-by-term with $x = q^n$.  This identification is unconditional given
the bar Euler product formula for $Y(\frakg_{K3})$.  Verified for all
$25$ classes to $q$-degree~$12$.

\medskip\noindent
\textbf{(III) Niemeier Yangian landscape and Umbral groups.}
At each of the $23$ Niemeier lattices $N$ with root system $R(N)$
\textup{(}excluding Leech\textup{)}, the K3 Yangian parameters
$h_i$ specialise to the root structure of~$N$.  The Umbral
group~$G_N$ (Cheng--Duncan--Harvey,
\texttt{arXiv:1204.2779}; proved by Duncan--Griffin--Ono,
\texttt{arXiv:1503.01472}) acts on the specialised parameters.
Universal properties across the landscape:
\begin{center}
\small
\renewcommand{\arraystretch}{1.2}
\begin{tabular}{lcc}
Property & Value & Moduli dependence \\
\hline
Bar Euler product & $\eta(\tau)^{24}$ & independent \\
$\kappa_{\mathrm{ch}}$ (lattice VOA) & $24$ & independent \\
$\kappa_{\mathrm{ch}}$ (K3 $\sigma$-model) & $2$ & independent \\
Shadow class (lattice VOA) & $G$ & independent \\
Shadow class (K3 $\sigma$-model) & $M$ & independent \\
Umbral group $G_N$ & $M_{24}$ to $\bZ_2$ & varies
\end{tabular}
\end{center}
The Umbral group order ranges from $|M_{24}| = 244{,}823{,}040$ at
$A_1^{24}$ to $|\bZ_2| = 2$ at most root systems.  Intermediate cases
include $2.M_{12}$ (order $190{,}080$) at $A_2^{12}$ and $3.S_6$
(order $2{,}160$) at $D_4^6$.

\medskip\noindent
\textbf{(III-a) Umbral mock modular forms from the shadow tower.}
At each Niemeier point~$N$ with Coxeter number~$h = \ell$
(the lambency), the K3 Yangian $Y(\frakg_{K3}(N))$ produces the
Umbral mock modular form~$h_N(\tau)$ through the shadow tower
mechanism:
\begin{equation}\label{eq:umbral-uniform-construction}
Y(\frakg_{K3}(N))
\;\xrightarrow{\text{shadow tower}}\;
(\kappa_{\mathrm{ch}},\, S_k)
\;\xrightarrow{\text{polar}}\;
h\big|_{\text{polar}} = -\kappa_{\mathrm{ch}} = -2
\;\xrightarrow{\text{Rademacher}}\;
h_N(\tau)
\;\xrightarrow{\theta\text{-decomp}}\;
\{H_r^{(\ell)}\}.
\end{equation}
The construction is \emph{uniform}: the mechanism is identical for
all~$23$ cases.  The polar coefficient $h_N\big|_{q^{-r^2/(4\ell)}} =
-\kappa_{\mathrm{ch}} = -2$ is universal (topological invariant,
moduli-independent), the shadow class is~$M$ (K3 sigma model, universal),
and the Rademacher sum converges at weight~$1/2$.  What varies is the
lambency $\ell$ (from~$2$ to~$46$), the modular level~$4\ell$, the
number of mock modular components~$H_r^{(\ell)}$, and the Umbral
group~$G_N$.

The parameter specialisation at each Niemeier point uses the Coxeter
exponents~$m_1, \ldots, m_r$ of the root system components.  For a
component of type~$\frakg$ with Coxeter number~$h(\frakg)$, the
Yangian parameters are $h_k = m_k - h(\frakg)/2$, centred so that
$\sum_k h_k = 0$ per component.  For the full~$24$ parameters,
$\sum_{i=1}^{24} h_i = 0$ (CY$_2$ constraint).  Verified for all~$23$
root systems.

The first few half-multiplicities (= $G_N$ representation dimensions):
\begin{center}
\small
\renewcommand{\arraystretch}{1.15}
\begin{tabular}{llcccccc}
$R(N)$ & $G_N$ & $\ell$ & $a_{-1}$ & $a_1$ & $a_2$ & $a_3$ & $a_4$ \\
\hline
$A_1^{24}$ & $M_{24}$ & $2$ & $-1$ & $45$ & $231$ & $770$ & $2277$ \\
$A_2^{12}$ & $2.M_{12}$ & $3$ & $-1$ & $16$ & $55$ & $144$ & $330$ \\
$A_3^{8}$  & $2.\mathrm{AGL}_3(2)$ & $4$ & $-1$ & $7$ & $21$ & $48$ & $99$ \\
$A_4^{6}$  & $\mathrm{GL}_2(5)/\bZ_2$ & $5$ & $-1$ & $4$ & $10$ & $20$ & $40$ \\
$D_4^{6}$  & $3.S_6$ & $6$ & $-1$ & $3$ & $7$ & $14$ & $27$ \\
$E_8^{3}$  & $S_3$ & $30$ & $-1$ & $1$ & & & \\
$D_{24}$   & $\bZ_2$ & $46$ & $-1$ & & & &
\end{tabular}
\end{center}
The polar coefficient~$a_{-1} = -1$ (giving $A_{-1} = \kappa_{\mathrm{ch}}
\cdot (-1) = -2$) is universal.  Coefficients at~$a_n$ for $n \geq 1$
are positive and decrease with lambency (reflecting smaller Umbral groups).
All~$23$ coefficient tables verified against
Cheng--Duncan--Harvey (\texttt{arXiv:1204.2779}).

\medskip\noindent
\textbf{(IV) Surfing group preservation under Yangian deformation.}
By the Gaberdiel--Hohenegger--Volpato theorem
(\texttt{arXiv:1008.3778}, \texttt{arXiv:1106.4315}), for every K3
sigma model~$M$, the symmetry group $G(M) < M_{24}$, and the union
$\bigcup_M G(M)$ generates $M_{24}$.  The Yangian deformation
PRESERVES this ``symmetry surfing'' structure: since $Y(\frakg_{K3})$
depends on $h_i$ only through the symmetric structure function
$g_{K3}(z) = \prod (z - h_i)/(z + h_i)$, any permutation of the
$h_i$ preserving $\sum h_i = 0$ is an automorphism of the abstract
algebra.  Thus $G(M)$ acts on $\Rep(Y(\frakg_{K3}(M)))$ at each point
$M$ in K3 moduli space.  The deformation parameter $\varepsilon$
(Yangian spectral shift) is a scalar and hence
$G(M)$-invariant.  The deformation neither enlarges nor reduces
$G(M)$: it preserves the surfing group exactly.

For the non-abelian case ($\frakg \neq \mathrm{gl}_1$), preservation
of $G(M)$ also requires $G(M) \subset \Aut(\omega_{ij})$,
where $\omega_{ij}$ is the Mukai pairing.  At the $A_1^{24}$ point
$\omega \propto I$, so $\Aut(\omega) = S_{24}$ and the condition
is vacuous.
\end{conjecture}

\noindent\textit{Verification}: $50$ tests in
\texttt{mathieu\_yangian\_deeper.py} covering generator permutation at
the $A_1^{24}$ point ($10$ tests), cyclotomic identity for all~$15$
cycle lengths and bar Euler = eta product for all $25$ Frame shapes
($12$ tests), Niemeier Yangian landscape with Umbral group orders
($11$ tests), surfing group preservation at Kummer, Fermat quartic,
$D_4^6$, and $E_8^3$ points ($10$ tests), and $7$ cross-verification
tests against the parent engine.
A further $80$ tests in \texttt{umbral\_23\_niemeier\_yangian.py}
covering parameter specialisation at all~$23$ Niemeier points ($15$~tests),
Umbral Frame shape eta products ($14$~tests), mock modular forms and
the shadow tower mechanism ($14$~tests), the uniform construction
$Y \to h_N$ ($8$~tests), Umbral group structure ($12$~tests),
landscape summary ($7$~tests), and $11$~multi-path cross-verification
tests, each carrying an explicit two-source verification trail.  A cross-check against the parent
engine~\texttt{mathieu\_moonshine\_yangian.py} corrected the $n = 7$
half-multiplicity from $30{,}492$ to $30{,}498$ ($= A_7/2 = 60{,}996/2$).

\subsection{The Kummer point and the Fermat quartic}
\label{subsec:k3-kummer-fermat-enhancement}

The Kummer point $K3 = T^4/\bZ_2$ and the Fermat quartic
$\{x_1^4 + x_2^4 + x_3^4 + x_4^4 = 0\} \subset \bC\bP^3$
are two explicitly constructed enhanced symmetry points that
illustrate distinct enhancement mechanisms.

\begin{remark}[The Kummer point: $\mathfrak{so}_4^{\oplus 4}$ at level~$1$]
\label{rem:k3-kummer-enhancement}
At the Kummer point (Steps~1--4 of
Proposition~\ref{prop:kummer-orbifold}):
\begin{enumerate}[label=(\roman*)]
\item $8$ untwisted Heisenberg generators from
 $H^{\mathrm{even}}(T^4)$, contributing $\kappa_{\mathrm{ch}} = 0$.
\item $16$ twisted sectors at the $A_1$ fixed points, enhancing
 to $\hat{\mathfrak{so}}_4^{\oplus 4}$ at level~$1$
 (Nahm--Wendland), contributing
 $\kappa_{\mathrm{ch}} = 4 \times \tfrac{1}{2} = 2$.
\item Total: $\kappa_{\mathrm{ch}}(\cA^{\mathrm{orb}}) = 0 + 2 = 2 = \chi(\cO_{K3})$.
\end{enumerate}
The enhanced subalgebra
$\hat{\mathfrak{so}}_4 \cong \hat{\mathfrak{su}}_2 \oplus \hat{\mathfrak{su}}_2$
has $c = 1 + 1 = 2$ per copy, central charge $4 \times 2 = 8$ total.
Each $\hat{\mathfrak{su}}_2$ at level~$1$ has shadow class~$L$ (depth~$3$,
$S_3 \neq 0$, $S_4 = 0$).  The full K3 sigma model at the Kummer point
remains class~$M$ (Computation~\ref{comp:shadow-tower-k3}).
\end{remark}

\begin{remark}[The Fermat quartic: Gepner model and discrete symmetry]
\label{rem:k3-fermat-quartic}
The Fermat quartic K3 is described by the Gepner model $(2)^4$:
four copies of the $\cN = 2$ minimal model at $k = 2$
($c = 3/2$ each), with GSO projection and $\bZ_4$ orbifold
(Remark~\ref{rem:gepner-model-k3}).

The discrete symmetry group at the Gepner point is
$(\bZ/4\bZ)^4 \rtimes S_4$ (order $6144$), embedding into the
Conway group $\mathrm{Co}_1$.  This large discrete symmetry
is the analogue of the $M_{24}$ action on the elliptic genus
(Mathieu moonshine, Theorem~\ref{thm:mathieu-moonshine}).

The Gepner tensor product and the physical sigma model give
\emph{different} modular characteristics:
\[
 \kappa_{\mathrm{Gepner}} = 4 \times \tfrac{3}{4} = 3
 \neq 2 = \kappa_{\mathrm{ch}}(\cA_{K3}).
\]
The discrepancy illustrates the trichotomy of
Proposition~\ref{prop:k3-trichotomy}: different constructions
extract different chiral algebras from the same geometry.
The Gepner tensor product uses the $\cN = 2$ Virasoro structure
of each factor ($\kappa_{\mathrm{ch}} = c/2$ for each), while the
physical sigma model uses the full $\cN = 4$ SCA which constrains
$\kappa_{\mathrm{ch}}$ to the topological value $\chi(\cO_{K3}) = 2$.
\end{remark}

\subsection{Wall-crossing structure between enhancement points}
\label{subsec:k3-enhancement-wall-crossing}

\begin{remark}[Wall-crossing preserves $\kappa_{\mathrm{ch}}$]
\label{rem:k3-wall-crossing-kappa}
Moving between enhanced symmetry points in $\cM_{K3}$ crosses walls
in the Bridgeland stability manifold
(Section~\ref{sec:k3e-wall-crossing}).  These wall-crossings
correspond to flops, birational transformations that are
derived autoequivalences of $D^b(\Coh(S))$. By
\[
 \kappa_{\mathrm{ch}}(A_X) = \kappa_{\mathrm{ch}}(A_{X^+})
 \quad \text{(flop $X \dashrightarrow X^+$ preserves $\kappa_{\mathrm{ch}}$)}.
\]
This is \emph{distinct} from the Koszul dual, which satisfies
$\kappa_{\mathrm{ch}}(A) + \kappa_{\mathrm{ch}}(A^!) = \rho_K$
(the family-dependent Koszul conductor), and which exchanges
algebra and coalgebra rather than chambers in the stability manifold.

The wall structure is encoded in the BKM root system of
$\mathfrak{g}_{\Delta_5}$
(Proposition~\ref{prop:k3e-walls-roots}): real roots produce
flop walls, lightlike roots produce null walls of multiplicity
$f(0) = 10$, and imaginary roots produce BPS walls with
multiplicity $|f(D)|$.  The Weyl group $W^{(2)}(\Lambda^{2,1}_{II})$
permutes chambers, and the denominator identity is the signed
chamber sum (Theorem~\ref{thm:k3e-denominator}).

The accumulation of infinitely many walls as one moves from
the large-volume chamber to the deep interior of the stability
manifold reflects the infinite root system of $\mathfrak{g}_{\Delta_5}$
and the class-$M$ shadow depth of the K3 sigma model.
\end{remark}


\subsection{Shadow class variation over K3 moduli}
\label{subsec:shadow-class-moduli-variation}

While $\kappa_{\mathrm{ch}}(K3) = 2$ is a topological invariant, constant
across the entire moduli space $\cM_{K3}$
(Theorem~\ref{thm:k3-kappa-moduli-invariance}), the
\emph{shadow class} is a finer invariant that varies as the K3 surface
moves through moduli.  The shadow class depends on the $A_\infty$
structure of the chiral algebra (specifically on the higher products
$m_3, m_4, \dotsc$), which detects the OPE structure that
$\kappa_{\mathrm{ch}}$ alone cannot see.

\begin{proposition}[K3 shadow class landscape]
\label{prop:k3-shadow-class-landscape}
\ClaimStatusProvedHere
Let $S$ be a K3 surface.  The shadow class of the chiral algebra
$\Phi_2(D^b(\Coh(S)))$ depends on the position of $S$ in $\cM_{K3}$
as follows.
\begin{enumerate}[label=(\roman*)]
\item \emph{Generic point} ($\rho = 1$, no enhanced symmetry):
 the chiral algebra is the Mukai Heisenberg $H_{\mathrm{Muk}}$
 (rank~$24$, signature~$(4,20)$).  The OPE is free-field:
 \[
  J^i(z)\,J^j(w) \sim \frac{\omega^{ij}}{(z-w)^2},
  \qquad i,j = 1,\dotsc,24.
 \]
 Shadow tower: $S_3 = S_4 = 0$, $\Delta = 0$.
 \textbf{Class~$G$, depth~$2$}.

\item \emph{ADE singularity point} ($\rho > 1$, enhanced $\hat{\fg}_1$
 subalgebra): the chiral algebra acquires nonabelian currents
 $J^a(z)$ with OPE
 \[
  J^a(z)\,J^b(w) \sim \frac{k\,\delta^{ab}}{(z-w)^2}
  + \frac{if^{ab}{}_c\,J^c(w)}{z - w},
 \]
 where $f^{ab}{}_c$ are the structure constants of $\fg$.
 The cubic OPE produces $S_3 \neq 0$; at level~$1$ the Sugawara
 structure gives $S_4 = 0$, hence $\Delta = 0$.
 Enhanced subalgebra: \textbf{class~$L$, depth~$3$}.
 Full $\cN = 4$ sigma model: \textbf{class~$M$}
 ($S_3 = 2$, $S_4 = 5/156$, $\Delta = 20/39 \neq 0$).

\item \emph{Kummer point} ($T^4/\bZ_2$, $\rho = 16$):
 $16$ twisted sectors at $A_1$ fixed points enhance to
 $\hat{\mathfrak{so}}_4^{\oplus 4}$ at level~$1$ (Nahm--Wendland).
 Each $\hat{\mathfrak{su}}_2$ at level~$1$ has $S_3 \neq 0$,
 $S_4 = 0$.  Enhanced subalgebra: \textbf{class~$L$}.
 Full sigma model: \textbf{class~$M$}.

\item \emph{Fermat quartic / Gepner point} ($\rho = 20$):
 the Gepner model $(2)^4/\bZ_4$ is a rational $\cN = (4,4)$ SCFT.
 As a standalone RCFT: \textbf{class~$L$} (finite shadow depth).
 Full sigma model: \textbf{class~$M$}.

\item \emph{Niemeier maximal points} ($24$ lattices, $\rho = 20$):
 the lattice VOA $V_N$ is class~$G$ (free-field, $S_3 = S_4 = 0$);
 the enhanced $\hat{\fg}_1$ subalgebra (for root system $R(N) \neq
 \varnothing$) is class~$L$.  The Leech lattice ($R = \varnothing$)
 is class~$G$.
\end{enumerate}
\end{proposition}

\begin{proof}
At the generic point, $\Phi_2(D^b(\Coh(S)))$ is the rank-$24$ Mukai
Heisenberg $H_{\mathrm{Muk}}$ by Theorem~\ref{thm:cy-a-d2}.
The Heisenberg OPE is free-field, so $S_3 = S_4 = 0$ and the shadow
class is~$G$ by the single-line dichotomy theorem (Vol~I,
Theorem~\ref{thm:shadow-dichotomy}).

At an ADE enhancement point, the chiral algebra acquires the affine
Kac--Moody subalgebra $\hat{\fg}_1$.  At level~$1$ for simply-laced
$\fg$: $c(\hat{\fg}_1) = \mathrm{rank}(\fg)$, and the structure constants
produce $S_3 \neq 0$ (cubic bar obstruction from the $J^aJ^bJ^c$
three-point function).  The quartic invariant $S_4 = 0$ at level~$1$
because the Sugawara tensor $T = (2(k + h^\vee))^{-1} {:}J^aJ^a{:}$ at
$k = 1$ has the same singular structure as the free-field stress tensor.
Hence $\Delta = 0$ and the enhanced subalgebra is class~$L$, depth~$3$.

The full $\cN = 4$ sigma model at $c = 6$ has the nonlinear $G^a$-$G^b$
OPE involving composite operators $J^aJ^b$.  This produces
$S_3 = 2$, $S_4 = 5/156$, hence $\Delta = 8 \cdot 2 \cdot 5/156 =
20/39 \neq 0$, giving class~$M$ (infinite shadow depth) by
Computation~\ref{comp:shadow-tower-k3}.

The Kummer, Fermat, and Niemeier cases follow from the same analysis
applied to the specific enhanced algebras at those moduli points.
\end{proof}

\noindent\textit{Verification}: $88$ tests in
\texttt{shadow\_class\_moduli\_variation.py}:
$20$+ moduli points ($8$ ADE types $A_1$--$A_8$, $5$ types $D_4$--$D_8$,
$3$ types $E_6$--$E_8$, generic, Kummer, Fermat, $4$ Niemeier);
$\kappa_{\mathrm{ch}} = 2$ at all points;
shadow class upper-semicontinuity at all transitions;
$15$ cross-checks between tower computation, point data,
Hodge theory, and stratification paths.

\begin{conjecture}[Upper-semicontinuity of shadow class]
\label{conj:shadow-usc}
\ClaimStatusConjectured
Let $(\cC_t)_{t \in B}$ be a flat family of $\CY_d$ categories over
a connected base~$B$.  For each $t \in B$, let $A_t = \Phi_d(\cC_t)$
be the chiral algebra (assuming $\CY\text{-}A_d$ applies) and let
$\mathrm{class}(t) \in \{G, L, C, M\}$ be the shadow class of~$A_t$.
Then the function $t \mapsto \mathrm{class}(t)$ is
upper-semicontinuous in the ordering $G < L < C < M$.

Equivalently: the locus $\{t \in B : \mathrm{class}(t) \geq X\}$
is closed for each $X \in \{G, L, C, M\}$, and the shadow depth
function $t \mapsto \mathrm{depth}(t)$ is lower-semicontinuous.
\end{conjecture}

\begin{remark}[Evidence for upper-semicontinuity]
\label{rem:shadow-usc-evidence}
At $d = 2$ (K3), Conjecture~\ref{conj:shadow-usc} is verified by
Proposition~\ref{prop:k3-shadow-class-landscape}: the generic class
is~$G$ (depth~$2$), and specialization to any enhanced symmetry point
produces class~$L$ (subalgebra) or~$M$ (sigma model), both~$\geq G$.

The heuristic: specialization introduces new OPE singularities (from
enhanced symmetry, degeneration, orbifold sectors, etc.)\ which can
only \emph{increase} the shadow depth.  Smoothing (moving to generic
position) can kill higher-order OPE terms, decreasing shadow depth.
This is analogous to the upper-semicontinuity of fiber dimension
in algebraic geometry.

At $d = 3$ (quintic family $X_\psi$, at large-complex-structure,
conifold, and Gepner degenerations), all studied points have conjectural
class~$M$ (conditional on CY-A$_3$): GV invariants grow
exponentially (Huang--Klemm--Quackenbush), forcing infinite shadow
depth.  Upper-semicontinuity holds trivially ($M \leq M$) but does
not provide a nontrivial test.  \textbf{Open question}: does there
exist a CY$_3$ family with generic shadow class $\neq M$?
\end{remark}

\begin{remark}[The conifold transition does not change shadow class]
\label{rem:conifold-shadow-class}
At the conifold point of the quintic family ($\psi^5 = 1$,
the conifold degeneration with vanishing $S^3$), a massless
hypermultiplet appears whose OPE is logarithmic, producing class~$M$.
The conifold transition $X_\psi \to X_0 \to X'$ changes the Hodge
numbers ($h^{1,1}: 1 \to 2$, $h^{2,1}: 101 \to 100$), the Euler
characteristic ($\chi: {-}200 \to {-}196$), and the shadow tower
\emph{coefficients}~$S_k$, but does \emph{not} change the shadow
class: it remains~$M$ on both sides and at the singular fiber.
This is because class~$M$ (infinite shadow depth) is a
\emph{robust} condition: perturbations of the BPS spectrum
change the shadow coefficients but not the qualitative
class.

All CY$_3$ shadow class statements are \emph{conjectural},
conditional on CY-A$_3$.
\end{remark}

\begin{remark}[The K3 moduli stratification]
\label{rem:k3-moduli-stratification}
The shadow class provides a stratification of the K3 moduli space:
\[
 \cM_{K3} = \cM_G \;\sqcup\; \cM_L \;\sqcup\; \cM_M
\]
where $\cM_G$ is the generic locus (class~$G$, full measure, open dense),
$\cM_L$ is the enhanced subalgebra locus (class~$L$, measure zero,
countable union of subvarieties of codimension~$\geq 1$), and
$\cM_M$ is the full sigma model at all non-generic points
(class~$M$, the complement of~$\cM_G$, also measure zero).
There is no class~$C$ stratum for K3: the jump from the generic
point is either $G \to L$ (subalgebra) or $G \to M$ (sigma model),
with no intermediate passage through~$C$.
\end{remark}
